% !TeX root = ../../Book3.tex
% 纲要标题：## 第17章　Gauss–Bonnet–陈定理与指标理论
% 纲要位置：几何学计划纲要.md，第 2782 行。
% 正文状态：初稿；尚未正式审核。

% 纲要内容开始
% 
% > 本章目标与证明范围：高维 Euler 公式保留完整证明。自旋与一般指标理论区分“构造可以证明”“局部公式需要更多分析”“一般指标定理外引”三个层次，不能把 Euler 特例扩大成一般证明。
% 
% 纲要内容结束

\chapter{Gauss–Bonnet–陈定理与指标理论}
\label{g3:ch:17}
\BookChapterNumber{g3:ch:17}

\begin{chapterabstract}
	曲率的局部表达可以在积分后成为拓扑整数。本章从 Pfaffian 构造 Euler 形式，证明其闭性与联络无关性，再在单位球丛上构造归一化转移形式。去掉向量场的孤立零点后，Stokes 定理把曲率积分化为边界映射度；结合 Poincaré--Hopf 定理，得到任意闭定向偶维黎曼流形的 Gauss--Bonnet--陈公式。选读部分构造 Clifford 模与自旋 Dirac 算子，推导 Lichnerowicz 公式，并说明正标量曲率对调和自旋子的限制。最后区分解析指标与符号的拓扑指标，准确引入一般指标定理，比较 de Rham、signature 和 spin 三种模型。
\end{chapterabstract}

\begin{learninggoals}
	\begin{itemize}
		\item 从外代数定义 Pfaffian，核对换架、直和与定向反转的符号。
		\item 计算 Euler 形式及联络路径的转移形式，恢复曲面中的 \(K\,\mathrm{vol}/(2\pi)\)。
		\item 在单位球丛上构造角形式，证明其纤维积分为一，并追踪去球后的边界定向。
		\item 把局部零点指标与曲率积分联系起来，完成高维 Euler 积分公式的证明。
		\item 构造自旋子丛及其 Dirac 算子，推导曲率收缩中的 \(\operatorname{Scal}/4\)。
		\item 识别椭圆符号的紧支 \(K\) 类，分清指标定理的外部输入与三个模型中可直接计算的解析指标。
	\end{itemize}
\end{learninggoals}

半径为 \(a\) 的圆球面，其 Gaussian 曲率为 \(a^{-2}\)，面积为 \(4\pi a^2\)，所以总曲率始终是 \(4\pi\)。改变半径改变了每一点的曲率和面积，却没有改变积分。曲面 Gauss--Bonnet 定理说明，这个不变量等于 \(2\pi\chi(S^2)\)。到了四维及更高偶数维，曲率不再由一个函数概括，首先要找出究竟应积分哪个曲率组合。

\ProseParagraphBreak

Pfaffian 把曲率二形式按定向配对，给出所需的 Euler 形式。其积分为何与向量场零点的总指标相同，则需要一个几何转换。陈省身的内蕴证明把问题提升到单位切球丛：每一个单位向量都在拉回切丛中给出一个典范非零向量，因此原来可能无法消除的零点障碍，在球丛上变成了一个边界积分。\footnote{陈省身（Shiing-Shen Chern，1911--2004），美籍华人数学家，研究整体微分几何、示性类与联络理论。}本章沿这条路线构造转移形式并计算全部归一化常数，原始论证见\cite[第2--3节，第749--752页]{Chern1944GaussBonnet}。

\ProseParagraphBreak

\chapref{g3:ch:15}还给出另一种计算 Euler 示性数的方法：把各次数调和形式的维数交错相加。令 \(D=\mathrm d+\delta\)，则 \(D^2=\Delta\)，而 \(D\) 在偶次与奇次形式之间交换。每个正特征值的两边模态可以配对，最后留下零模维数之差。这使曲率积分、上同调和椭圆算子的核出现在同一问题中，也自然引出指标观点。

\ProseParagraphBreak

前两节构成本章主干，使用\chapref{g3:ch:03}的曲率、微分形式与 Stokes 定理，以及第二卷第16.3节的局部指标和第17.4.3小节已证的 Poincaré--Hopf 定理。后两节为选读，需要\chapref{g3:ch:13}的椭圆 Fredholm 理论以及\chapref{g3:ch:15}、\chapref{g3:ch:16}的 Hodge、热核工具；涉及拓扑 \(K\) 理论时，就近恢复第二卷第24章的符号类和推前构造。一般指标等式采用明确的外部定理，定位见\cite[第6节，定理(6.7)，第520页]{AtiyahSinger1968I}及\cite[第5--6节]{AtiyahSinger1968III}。

% !TeX root = ../../Book3.tex
% 纲要标题：### 17.1 Euler 形式的构造
% 纲要位置：几何学计划纲要.md，第 2786 行。
% 正文状态：初稿；尚未正式审核。

% 纲要内容开始
% 
% * Pfaffian 的线性代数
% * 曲率形式和归一化
% * 陈–Weil 闭性及联络无关性
% 
% 纲要内容结束

\section{Euler 形式的构造}
\label{g3:sec:1701}

曲面上的曲率积分使用 \(K\,\mathrm dA\)：曲率给出函数，定向面积形式把它变成可积分的二形式。在 \(2m\) 维，曲率本身是一个取值于斜对称矩阵的二形式。我们需要从这个矩阵中提取一个 \(2m\) 次形式，并使它在正交标架变化时不变、在定向反转时变号。Pfaffian 恰好具有这些性质。

\subsection{Pfaffian 的线性代数}
\label{g3:subsec:1701:01}
先把流形放在一边。设 \(V\) 为定向实 \(2m\) 维内积空间，\(m\ge1\)，并选取正向正交基及其对偶基 \(\varepsilon^1,\ldots,\varepsilon^{2m}\)。对斜对称实矩阵 \(X=(X_{ij})\)，令
\[
	\eta_X=\frac12\sum_{i,j=1}^{2m}X_{ij}\varepsilon^i\wedge\varepsilon^j.
\]
顶外幂只有一个方向，因此 \(\eta_X^m/m!\) 必为体积元的倍数。这个倍数称为 \(X\) 的\term[Pfaffian]{Pfaffian}[Pfaffian]。\footnote{Johann Friedrich Pfaff（1765--1825），德国数学家，研究微积分与一阶偏微分方程及其微分系统。}
\begin{definition}\label{g3:def:pfaffian}
	设 \(X\) 为实 \(2m\) 阶斜对称矩阵，\(m\ge1\)。定义 \(m\) 次齐次多项式
	\begin{equation}\label{g3:eq:pfaffian-definition}
		\operatorname{Pf}(X)
		=\frac1{2^m m!}\sum_{\sigma\in S_{2m}}\operatorname{sgn}(\sigma)
		X_{\sigma(1)\sigma(2)}\cdots X_{\sigma(2m-1)\sigma(2m)}.
	\end{equation}
	令 \(\varepsilon^i\) 为 \(\mathbb R^{2m}\) 的标准对偶基、\(\eta_X=\tfrac12\sum_{i,j}X_{ij}\varepsilon^i\wedge\varepsilon^j\)。等价地，\(\eta_X^m/m!=\operatorname{Pf}(X)\varepsilon^1\wedge\cdots\wedge\varepsilon^{2m}\)。空矩阵的 Pfaffian 约定为 \(1\)。
\end{definition}
\symidx{\(\operatorname{Pf}(X)\)：斜对称矩阵的 Pfaffian}

展开 \(\eta_X^m\) 时，重复的基指标产生零项；其余项由 \(1,\ldots,2m\) 的排列给出，正是式\eqref{g3:eq:pfaffian-definition}。分母中的 \(2^m\) 来自每一对指标的两个顺序，\(m!\) 来自各二形式因子的排列。

\begin{proposition}\label{g3:prop:pfaffian-algebra}
	设 \(X\) 为实 \(2m\) 阶斜对称矩阵，\(m\ge1\)，\(B\) 为任意实 \(2m\) 阶矩阵。则
	\begin{equation}\label{g3:eq:pfaffian-change-frame}
		\operatorname{Pf}(B^{\mathsf T}XB)=\det(B)\operatorname{Pf}(X),
		\qquad \operatorname{Pf}(X)^2=\det X.
	\end{equation}
	若 \(X=X_1\oplus X_2\)，且两个块均为偶数阶，按先第一块、后第二块定向，则
	\(\operatorname{Pf}(X)=\operatorname{Pf}(X_1)\operatorname{Pf}(X_2)\)。这些恒等式也适用于实交换代数中的矩阵系数。
\end{proposition}
\begin{proof}
	把 \(B\) 看作 \(V\) 上的线性映射，则 \(B^*\eta_X=\eta_{B^{\mathsf T}XB}\)，而 \(B^*\) 在顶外幂上乘 \(\det B\)。对定义中的顶次式拉回即得第一式，包括 \(B\) 不可逆的情形。分块时 \(\eta_X=\eta_{X_1}+\eta_{X_2}\)；能在各块同时达到顶次的项只有相应的幂次组合，其二项式系数被阶乘抵消，得到乘法公式。

	为证明平方公式，可逐次把实斜对称矩阵化为正交二阶块。若 \(X\ne0\)，对称非负算子 \(-X^2\) 有一个正特征值 \(a^2\)。取对应单位向量 \(v\)，令 \(w=-Xv/a\)，则 \(v,w\) 正交且 \(Xv=-aw\)、\(Xw=av\)。它们张成的平面及其正交补均在 \(X\) 下不变。重复这个过程，剩下的零空间维数为偶数，可以任意两两配对。因此存在正交矩阵 \(Q\)，使
	\[
		Q^{\mathsf T}XQ=\bigoplus_{j=1}^m
		\begin{pmatrix}0&a_j\\-a_j&0\end{pmatrix}.
	\]
	右侧的 Pfaffian 为 \(\prod_j a_j\)，行列式为 \(\prod_j a_j^2\)。换基式和 \((\det Q)^2=1\) 给出所需平方公式。最后，以上恒等式都是实系数多项式恒等式：对所有实矩阵成立便说明各单项式系数相等，因此可以把变量代入任意实交换代数。
\end{proof}

例如
\[
	\operatorname{Pf}\begin{pmatrix}0&a\\-a&0\end{pmatrix}=a,
	\qquad
	\operatorname{Pf}\begin{pmatrix}
		0&a&b&c\\-a&0&d&e\\-b&-d&0&f\\-c&-e&-f&0
	\end{pmatrix}=af-be+cd.
\]
平方公式说明 Pfaffian 保留了行列式丢失的定向信息：把一个基向量反号，Pfaffian 反号，行列式不变。随后把 \(a,b,\ldots\) 换成二形式时，乘法就是外乘；偶次形式彼此交换，所以上述多项式计算仍然适用。

\ProseParagraphBreak

改变联络时，需要把曲率多项式对某个方向求导。为此用 \(\operatorname{Pf}(X_1,\ldots,X_m)\) 表示 Pfaffian 的对称多重线性极化，具体定义为
\begin{equation}\label{g3:eq:pfaffian-polarization}
	\operatorname{Pf}(X_1,\ldots,X_m)
	=\left.\frac1{m!}\frac{\partial^m}{\partial t_1\cdots\partial t_m}
	\operatorname{Pf}\left(\sum_{r=1}^m t_rX_r\right)\right|_{t_1=\cdots=t_m=0}.
\end{equation}
它满足 \(\operatorname{Pf}(X,\ldots,X)=\operatorname{Pf}(X)\)。对矩阵值微分形式，先写成实矩阵与标量形式的有限和，再按参数顺序外乘标量形式，便得到同一多重线性表达。若两个参数的形式次数为 \(p,q\)，交换它们带符号 \((-1)^{pq}\)。特别地，一个一形式与若干二形式之间的换位都不产生符号。

\subsection{曲率形式和归一化}
\label{g3:subsec:1701:02}
设 \(E\to M\) 是光滑流形上的定向实秩 \(2m\) 向量丛，带光滑内积和度量联络 \(\nabla\)。局部选取正向正交标架 \(e=(e_1,\ldots,e_{2m})\)，以列指标记录被微分的向量：
\begin{equation}\label{g3:eq:euler-connection-curvature}
	\nabla e_j=\sum_i e_i\omega_{ij},\qquad
	\Omega=\mathrm d\omega+\omega\wedge\omega,\qquad
	\Omega_{ij}(X,Y)=\langle R^E(X,Y)e_j,e_i\rangle.
\end{equation}
度量相容性给出 \(\omega^{\mathsf T}=-\omega\)、\(\Omega^{\mathsf T}=-\Omega\)。这里 \(R^E(X,Y)=[\nabla_X,\nabla_Y]-\nabla_{[X,Y]}\)，与本卷的向量曲率次序相同。

\begin{definition}\label{g3:def:euler-form}
	设 \(E\to M\) 为光滑流形上的定向实秩 \(2m\ge2\) 内积丛，\(\nabla\) 为度量联络。定义其\term[Euler xingshi]{Euler 形式}[Euler form]为
	\begin{equation}\label{g3:eq:euler-form}
		e(E,\nabla)=e(\nabla)\coloneqq\frac1{(2\pi)^m}\operatorname{Pf}(\Omega)
		\in\Omega^{2m}(M;\mathbb R),
	\end{equation}
	其中在正向正交标架 \(e_j\) 中写 \(\nabla e_j=\sum_i e_i\omega_{ij}\)，\(\Omega=\mathrm d\omega+\omega\wedge\omega\)；\(\operatorname{Pf}\) 为 Pfaffian。
\end{definition}
\symidx{\(e(\nabla)\)：度量联络的 Euler 形式}

这个局部定义确实能粘合。若 \(e'=eQ\)、\(Q:M\supset U\to\operatorname{SO}(2m)\)，则
\[
	\omega'=Q^{-1}\omega Q+Q^{-1}\mathrm dQ,\qquad
	\Omega'=Q^{-1}\Omega Q=Q^{\mathsf T}\Omega Q.
\]
由\cref{g3:prop:pfaffian-algebra}，\(\operatorname{Pf}(\Omega')=\operatorname{Pf}(\Omega)\)。同样的计算说明，反转丛定向使 \(e(\nabla)\) 反号，拉回丛及拉回联络的 Euler 形式则为原形式的拉回。

\begin{example}\label{g3:exm:euler-surface-normalization}
	设 \((M^2,g)\) 为定向光滑黎曼曲面，\((e_1,e_2)\) 为正向正交标架，\((\theta^1,\theta^2)\) 为对偶标架。写 \(\nabla e_1=\alpha e_2\)，则 \(\omega_{21}=\alpha\)、\(\omega_{12}=-\alpha\)，从而 \(\Omega_{12}=-\mathrm d\alpha\)。另一方面，本卷的截面曲率为 \(K=\langle R(e_1,e_2)e_2,e_1\rangle\)，所以
	\[
		\Omega_{12}=K\,\theta^1\wedge\theta^2,
		\qquad e(\mathrm TM,\nabla)=\frac{K}{2\pi}\,\mathrm{vol}_g.
	\]
	因此 \(2\pi\) 的归一化恢复了曲面 Gauss--Bonnet 公式中的曲率积分。
\end{example}

\begin{proposition}\label{g3:prop:euler-constant-curvature}
	设 \((M^{2m},g)\) 为定向光滑黎曼流形，\(m\ge1\)，截面曲率恒为 \(\kappa\)。其 Levi-Civita 联络满足
	\begin{equation}\label{g3:eq:euler-space-form}
		e(\nabla)=\frac{(2m)!}{2^{2m}m!\pi^m}\,\kappa^m\mathrm{vol}_g.
	\end{equation}
\end{proposition}
\begin{proof}
	式\eqref{g3:eq:constant-curvature-tensor}给出 \(R(X,Y)Z=\kappa(\langle Z,Y\rangle X-\langle Z,X\rangle Y)\)，故 \(\Omega_{ij}=\kappa\theta^i\wedge\theta^j\)。代入式\eqref{g3:eq:pfaffian-definition}，每个排列项的外乘符号与其前面的 \(\operatorname{sgn}(\sigma)\) 抵消，全部 \((2m)!\) 项均为 \(\kappa^m\mathrm{vol}_g\)。再乘 \((2\pi)^{-m}/(2^m m!)\) 即得结论。
\end{proof}

例如四维常曲率情形为 \(e(\nabla)=3\kappa^2\mathrm{vol}_g/(4\pi^2)\)。一般四维流形的曲率形式未必具有 \(\Omega_{ij}=\kappa\theta^i\wedge\theta^j\) 的形式；这时应使用完整表达
\[
	e(\nabla)=\frac1{4\pi^2}
	\bigl(\Omega_{12}\wedge\Omega_{34}-\Omega_{13}\wedge\Omega_{24}
	+\Omega_{14}\wedge\Omega_{23}\bigr).
\]
三组配对共同把不同曲率平面的信息组合成顶次形式。

\subsection{陈–Weil 闭性及联络无关性}
\label{g3:subsec:1701:03}
Euler 形式由某个联络计算，拓扑数却不应随联络的选择改变。关键是把联络变化造成的差写成一个外微分。这是\term[Chen Weil fangfa]{陈--Weil 方法}[Chern--Weil theory]在 Euler 多项式上的具体实现。\footnote{André Abraham Weil（1906--1998），法国数学家，研究数论、代数几何与群论。}

\begin{lemma}\label{g3:lem:pfaffian-covariant-differential}
	设 \(E\to M\) 为定向实秩 \(2m\ge2\) 内积丛，\(\nabla\) 为光滑度量联络。对取值于斜对称自同态的齐次微分形式 \(A_r\in\Omega^{p_r}(M;\mathfrak{so}(E))\)，有
	\begin{equation}\label{g3:eq:pfaffian-covariant-differential}
		\mathrm d\operatorname{Pf}(A_1,\ldots,A_m)
		=\sum_{r=1}^m(-1)^{p_1+\cdots+p_{r-1}}
		\operatorname{Pf}(A_1,\ldots,\mathrm d_\nabla A_r,\ldots,A_m).
	\end{equation}
	其中局部写作 \(\mathrm d_\nabla A=\mathrm dA+\omega\wedge A-(-1)^{\deg A}A\wedge\omega\)。曲率满足 \(\mathrm d_\nabla\Omega=0\)。
\end{lemma}
\begin{proof}
	先解释式\eqref{g3:eq:pfaffian-covariant-differential}的两端为何都是全局形式。Pfaffian 对 \(Q\in\operatorname{SO}(2m)\) 的不变性经极化后仍成立，\(A_r\) 及 \(\mathrm d_\nabla A_r\) 在换架时都按 \(Q^{-1}(\cdot)Q\) 变换，因此两端与正交标架无关。

	固定 \(x\in M\)，可选局部正向正交标架使 \(\omega(x)=0\)：从任意框架出发，在中心为 \(x\) 的坐标 \(z\) 中用 \(Q(z)=\exp(-\sum_a z^a\omega_a(x))\) 换架即可。在这一点，所有外协变微分均成为普通外微分，极化表达是若干标量形式外乘的常系数线性组合，外微分的乘积法则正给出所述符号。由于 \(x\) 任意，等式处处成立。

	最后直接展开 \(\Omega=\mathrm d\omega+\omega\wedge\omega\)：
	\[
		\mathrm d\Omega=\mathrm d\omega\wedge\omega-\omega\wedge\mathrm d\omega,
		\qquad
		\omega\wedge\Omega-\Omega\wedge\omega
		=\omega\wedge\mathrm d\omega-\mathrm d\omega\wedge\omega.
	\]
	相加为零，即 Bianchi 恒等式。
\end{proof}

\begin{theorem}[Euler 形式的陈--Weil 性质]\label{g3:thm:euler-chern-weil}
	设 \(E\to M\) 为光滑流形上的定向实秩 \(2m\ge2\) 内积丛，\(\nabla^0,\nabla^1\) 为光滑度量联络。令
	\[
		A=\nabla^1-\nabla^0\in\Omega^1(M;\mathfrak{so}(E)),\qquad
		\nabla^t=\nabla^0+tA,\quad 0\le t\le1,
	\]
	并记 \(\Omega_t\) 为 \(\nabla^t\) 的曲率。则 \(e(\nabla^j)\) 为闭形式，且定义在 \(M\) 上的 \((2m-1)\) 次形式
	\begin{equation}\label{g3:eq:euler-transgression-form}
		\operatorname{TP}(\nabla^0,\nabla^1)
		=\frac{m}{(2\pi)^m}\int_0^1
		\operatorname{Pf}(A,\underbrace{\Omega_t,\ldots,\Omega_t}_{m-1\text{ 个}})\,\mathrm dt
	\end{equation}
	满足
	\begin{equation}\label{g3:eq:euler-transgression}
		\mathrm d\operatorname{TP}(\nabla^0,\nabla^1)
		=e(\nabla^1)-e(\nabla^0).
	\end{equation}
	因而 \([e(\nabla)]\in H_{\mathrm{dR}}^{2m}(M)\) 与度量联络的选择无关。
\end{theorem}
\begin{proof}
	曲率为二形式，各参数之间交换时无符号。\cref{g3:lem:pfaffian-covariant-differential}和 Bianchi 恒等式给出
	\[
		\mathrm d\operatorname{Pf}(\Omega)=m\operatorname{Pf}(\mathrm d_\nabla\Omega,\Omega,\ldots,\Omega)=0.
	\]
	因此在任意维数的底空间上，Euler 形式都闭。

	在同一局部标架中 \(\omega_t=\omega_0+tA\)。对曲率求导得
	\[
		\frac{\mathrm d}{\mathrm dt}\Omega_t
		=\mathrm dA+\omega_t\wedge A+A\wedge\omega_t
		=\mathrm d_{\nabla^t}A.
	\]
	极化后的乘积法则遂给出
	\[
		\begin{aligned}
			\frac{\mathrm d}{\mathrm dt}e(\nabla^t)
			&=\frac{m}{(2\pi)^m}\operatorname{Pf}(\mathrm d_{\nabla^t}A,\Omega_t,\ldots,\Omega_t)\\
			&=\mathrm d\left[\frac{m}{(2\pi)^m}\operatorname{Pf}(A,\Omega_t,\ldots,\Omega_t)\right].
		\end{aligned}
	\]
	第二步中所有含 \(\mathrm d_{\nabla^t}\Omega_t\) 的项都为零。\(A\) 是两个联络的差，按自同态值一形式变换；所以被积形式全局定义，而不是依赖某张标架的局部原函数。在有限区间上积分，与外微分交换，就得到式\eqref{g3:eq:euler-transgression}。
\end{proof}

形式 \(\operatorname{TP}\) 称为这两个联络之间的\term[zhuanyi xingshi]{转移形式}[transgression form]。若基空间闭且定向、维数为 \(2m\)，Stokes 定理立即给出 \(\int_Me(\nabla^1)=\int_Me(\nabla^0)\)。下一节将把同一公式用在球丛上：在那里有一个全局非零截面，可构造曲率退化的联络，使 Euler 形式变成零，留下一个可精确计算的边界形式。

\begin{proposition}\label{g3:prop:euler-form-basic-properties}
	设 \(E\to M\) 为定向实偶数秩光滑向量丛。令 \(e(\nabla)=(2\pi)^{-m}\operatorname{Pf}(\Omega_\nabla)\)，其中 \(2m=\operatorname{rank}E\)。此 de Rham 类 \([e(\nabla)]\) 也不依赖内积。对定向偶数秩内积丛 \(E_1,E_2\) 及其度量联络，有
	\[
		e(\nabla^1\oplus\nabla^2)=e(\nabla^1)\wedge e(\nabla^2),
	\]
	其中直和采用先 \(E_1\)、后 \(E_2\) 的定向。若一个度量联络有处处非零的平行截面，则其 Euler 形式恒为零。
\end{proposition}
\begin{proof}
	对两个内积 \(h_0,h_1\)，存在唯一 \(h_0\) 对称正定丛自同态 \(B\)，使 \(h_1(v,w)=h_0(Bv,w)\)。令 \(S=B^{-1/2}\)，便有 \(h_1(Sv,Sw)=h_0(v,w)\)，且 \(S\) 保向。正定平方根光滑：在正定矩阵处，平方映射的导数 \(H\mapsto CH+HC\) 在对称矩阵空间可逆，因为在 \(C\) 的正交特征基中各分量乘以两个正特征值之和；局部反函数定理与正定平方根的唯一性给出光滑性。将一个 \(h_1\) 度量联络由 \(S\) 运回，得到 \(h_0\) 度量联络；它的 Euler 形式与原联络相同，再应用\cref{g3:thm:euler-chern-weil}即可。

	直和曲率矩阵为分块对角矩阵，Pfaffian 的乘法公式给出第二个结论。若 \(u\) 是处处非零的平行截面，逐连通分支把它归一化，并选以 \(u\) 为首向量的局部正交标架。由 \(R^Eu=0\)，曲率矩阵的首列及首行均为零，每个 Pfaffian 单项式都含其中一个因子，故全部为零。
\end{proof}

曲率形式的闭性来自 Bianchi 恒等式，联络无关性来自路径积分。两者都能在任意维数的底空间上使用；把所得顶次形式积分后为何得到整数，下一节将通过向量场零点的映射度来解释。
% 纲要条目排为编号小节标题，并进入目录与书签；说明段仍保留于纲要注释中。

% !TeX root = ../../Book3.tex
% 纲要标题：### 17.2 Gauss–Bonnet–陈定理的转移证明
% 纲要位置：几何学计划纲要.md，第 2792 行。
% 正文状态：初稿；尚未正式审核。

% 纲要内容开始
% 
% * 拉回丛的典范截面
% * 转移形式及其外微分
% * 纤维积分的归一化
% * 孤立零点附近的边界积分
% * Poincaré–Hopf 与整体 Euler 公式
% 
% 纲要内容结束

\section{Gauss–Bonnet–陈定理的转移证明}
\label{g3:sec:1702}

Euler 形式已经由曲率构造出来，但其积分为什么是一个拓扑整数，还需要把曲率与可计数的对象联系起来。第二卷第22.3.3小节曾通过 Euler 类的曲率代表给出这一比较。本节给出另一种完整证明：把 Euler 形式拉到单位球丛，利用那里的典范非零截面构造原形式；再用一个向量场把原形式拉回穿孔流形，孔洞的边界积分便读出零点指标。球丛转移与零点计数相结合的思路来自陈省身的内蕴证明\cite[第2--3节，第749--752页]{Chern1944GaussBonnet}。这里以两联络之间的转移公式组织计算，并逐项核对纤维上的常数与边界方向。

\subsection{拉回丛的典范截面}
\label{g3:subsec:1702:01}
先把切丛暂时换成一般的定向 Euclid 丛。设 \(E\to B\) 为光滑无边界流形上的定向实秩 \(2m\) 丛，\(m\ge1\)，带光滑纤维内积及相容联络 \(\nabla\)。其单位球丛是
\[
	S(E)=\{(x,v):x\in B,\ v\in E_x,\ |v|=1\},
	\qquad \pi:S(E)\longrightarrow B.
\]
正向正交平凡化把它写成 \(U\times S^{2m-1}\)；换架属于 \(\operatorname{SO}(2m)\)，因而保持球面方向。每根纤维 \(S(E_x)\) 都采用有向空间 \(E_x\) 中单位球的外法向优先边界取向。

\ProseParagraphBreak

拉回丛 \(\pi^*E\to S(E)\) 的纤维在 \((x,v)\) 处是 \(E_x\)，所以它有一个无须选择的单位截面
\begin{equation}\label{g3:eq:sphere-tautological-section}
	u(x,v)=v.
\end{equation}
这就是球丛上的\term[qiuceng dianfan danwei jiemian]{典范单位截面}[tautological unit section]。它不等于底丛 \(E\) 上的非零截面：每个底点上所有可能方向都仍保留在球丛中。记
\[
	D=\pi^*\nabla,\qquad w=Du\in\Omega^1(S(E);\pi^*E).
\]
度量相容性给出 \(\langle u,w\rangle=\frac12\mathrm d|u|^2=0\)。在正交标架中，\(u\) 是列向量，\(w=\mathrm du+\omega u\)；转置只作用于纤维指标。

\begin{lemma}\label{g3:lem:tautological-parallel-connection}
	设 \(E\to B\) 为光滑无边界流形上的定向实秩 \(2m\ge2\) Euclid 丛，\(\nabla\) 为相容联络。在其单位球丛上取式\eqref{g3:eq:sphere-tautological-section}的截面 \(u\)，令 \(D=\pi^*\nabla\)、\(w=Du\)。则
	\begin{equation}\label{g3:eq:tautological-parallel-connection}
		D^0=D+C,\qquad C=uw^{\mathsf T}-wu^{\mathsf T}
	\end{equation}
	是 \(\pi^*E\) 上的相容联络，满足 \(D^0u=0\)。它的 Euler 形式恒为零。
\end{lemma}
\begin{proof}
	这里的表达式内在地指
	\[
		C(Z)v=u\langle w(Z),v\rangle-w(Z)\langle u,v\rangle,
		\qquad Z\in\mathrm T S(E),\ v\in\pi^*E.
	\]
	所以 \(C\) 是全局定义的丛自同态值一形式。交换 \(v\) 与另一个纤维向量，可直接验证 \(C(Z)\) 斜对称；因此给相容联络 \(D\) 加上 \(C\) 后仍相容。又因 \(|u|=1\)、\(\langle u,w\rangle=0\)，有 \(Cu=-w\)，从而 \(D^0u=Du+Cu=0\)。

	记 \(D^0\) 的曲率为 \(\Omega^0\)。由 \((D^0)^2u=0\)，得到 \(\Omega^0u=0\)。在任一点附近，将 \(u\) 补成正向正交标架，其曲率矩阵的第一列为零，斜对称性使第一行也为零。Pfaffian 展开中的每个乘积都含指标 \(1\)，因而每项为零。故 \(e(D^0)=0\)。
\end{proof}

这个修改保留了正交结构，同时使一条单位方向平行。Pfaffian 必须使用全部 \(2m\) 个纤维方向；当曲率在其中一条方向上消失时，Euler 形式也就消失。这正是选择球丛的作用。

\subsection{转移形式及其外微分}
\label{g3:subsec:1702:02}
将上一节的联络转移公式用于 \(D^0\) 与 \(D\)。记
\[
	A=D-D^0=-C,\qquad D^t=D^0+tA,\qquad 0\le t\le1,
\]
并以 \(\Omega^t\) 表示 \(D^t\) 的曲率。定义 \(S(E)\) 上的 \(2m-1\) 形式
\begin{equation}\label{g3:eq:sphere-euler-transgression}
	\Psi_\nabla=-\frac{m}{(2\pi)^m}
	\int_0^1\operatorname{Pf}
	\bigl(A,\underbrace{\Omega^t,\ldots,\Omega^t}_{m-1\text{ 个}}\bigr)\,\mathrm dt.
\end{equation}
\symidx{\(\Psi_\nabla\)：单位球丛上的归一化 Euler 转移形式}
当 \(m=1\) 时没有曲率因子，积分式就是 \(-A_{12}/(2\pi)\)。一般情形中，一个一形式与 \(m-1\) 个二形式给出次数 \(2m-1\)。对称极化 \(\operatorname{Pf}(A,\Omega^t,\ldots,\Omega^t)\) 采用上一节的归一化。

\begin{proposition}\label{g3:prop:sphere-euler-transgression-differential}
	设 \(E\to B\) 为光滑无边界流形上的定向实秩 \(2m\ge2\) Euclid 丛，\(\nabla\) 为相容联络。式\eqref{g3:eq:sphere-euler-transgression}定义一个全局光滑实形式，并满足
	\begin{equation}\label{g3:eq:sphere-euler-transgression-differential}
		\mathrm d\Psi_\nabla=-\pi^*e(\nabla).
	\end{equation}
	此构造在保向等距丛同构及联络拉回下自然。
\end{proposition}
\begin{proof}
	\(D,D^0\) 都是同一 Euclid 丛上的全局相容联络，其差 \(A\) 是张量。换正交标架时，\(A\) 与 \(\Omega^t\) 同时共轭；Pfaffian 极化的不变性说明积分中的局部表达粘合成全局形式。对光滑被积形式在紧参数区间积分，也保持光滑性。

	由\cref{g3:thm:euler-chern-weil}及式\eqref{g3:eq:euler-transgression}，
	\[
		\mathrm d\Psi_\nabla
		=-\bigl(e(D)-e(D^0)\bigr)
		=-\pi^*e(\nabla).
	\]
	最后一步用了前一引理及曲率对拉回的自然性。所有构造只使用单位截面、联络、内积、外积和保向不变多项式；逐项拉回便得到自然性。
\end{proof}

等式前面的负号是有意选择的。稍后从底流形挖去小球时，新边界的外法向指向被挖去的球内；这会再产生一个负号。两者相消之后，正的局部指标才给 Euler 积分贡献正数。不过，仅有外微分等式还不能决定常数：还须证明 \(\Psi_\nabla\) 沿每根有向球面纤维的积分恰为 \(1\)。

\subsection{纤维积分的归一化}
\label{g3:subsec:1702:03}
固定 \(x\in B\)。沿纤维 \(S(E_x)\)，拉回联络 \(D\) 没有底方向的联络项；选 \(E_x\) 的一个正向正交基，就把它识别成固定向量空间 \(\mathbb R^{2m}\) 上的平凡联络 \(\mathrm d\)。因此这一限制上的全部计算只涉及单位球面的坐标 \(u_1,\ldots,u_{2m}\)，与底空间曲率无关。

\begin{theorem}\label{g3:thm:sphere-euler-angular-form}
	设 \(E\to B\) 为光滑无边界流形上的定向实秩 \(2m\ge2\) Euclid 丛，\(\nabla\) 为相容联络。式\eqref{g3:eq:sphere-euler-transgression}的形式满足
	\[
		\mathrm d\Psi_\nabla=-\pi^*e(\nabla),\qquad
		\int_{S(E_x)}\Psi_\nabla=1\quad(x\in B).
	\]
	更具体地，在每根纤维的正向正交坐标中，
	\begin{equation}\label{g3:eq:sphere-angular-normalization}
		\left.\Psi_\nabla\right|_{S(E_x)}
		=\frac{(m-1)!}{2\pi^m}\,\sigma_{2m-1},
		\qquad
		\sigma_{2m-1}=
		\sum_{i=1}^{2m}(-1)^{i-1}u_i\,
		\mathrm du_1\wedge\cdots\wedge\widehat{\mathrm du_i}
		\wedge\cdots\wedge\mathrm du_{2m}.
	\end{equation}
	这里 \(\sigma_{2m-1}\) 是外法向优先的标准球面体积形式。
\end{theorem}
\begin{proof}
	外微分公式已经证明。以下限制在一根纤维上；这时 \(w=\mathrm du\)，于是
	\[
		C=u\,\mathrm du^{\mathsf T}-\mathrm du\,u^{\mathsf T},
		\quad A=-C,\quad D^t=\mathrm d+(1-t)C,
		\quad \Omega^t=(1-t)\mathrm dC+(1-t)^2C\wedge C.
	\]
	整个表达在 \(\operatorname{SO}(2m)\) 的常值换架下不变。该群在球面上传递，所以只需在 \(u=e_1\) 处计算。在这里 \(\mathrm du_1=0\) 于球面的切空间，记 \(\theta_a=\mathrm du_a\)、\(2\le a\le2m\)。有
	\[
		C_{1a}=\theta_a,\quad C_{a1}=-\theta_a,\quad
		A_{1a}=-\theta_a,\quad A_{a1}=\theta_a.
	\]
	其余 \(C_{ab}\) 在此点为零，但它们的外微分不能省略。由完整公式求导，
	\[
		\mathrm dC=2\,\mathrm du\wedge\mathrm du^{\mathsf T},
		\qquad
		(\mathrm dC)_{ab}=2\theta_a\wedge\theta_b,
		\qquad
		(C\wedge C)_{ab}=-\theta_a\wedge\theta_b.
	\]
	因此
	\begin{equation}\label{g3:eq:angular-fiber-path-curvature}
		\Omega^t_{ab}=(1-t^2)\theta_a\wedge\theta_b
		\quad(a,b\ge2),\qquad
		\Omega^t_{1a}=\Omega^t_{a1}=0.
	\end{equation}

	现在展开极化 Pfaffian。因为只有一个奇次因子 \(A\)，将它越过二次曲率因子不产生符号，所以对称极化可写成
	\[
		\operatorname{Pf}(A,\Omega^t,\ldots,\Omega^t)
		=\frac1{2^m m!}\sum_{\tau\in S_{2m}}\operatorname{sgn}(\tau)
		A_{\tau(1)\tau(2)}
		\wedge\Omega^t_{\tau(3)\tau(4)}
		\wedge\cdots\wedge\Omega^t_{\tau(2m-1)\tau(2m)}.
	\]
	非零项必须让指标 \(1\) 出现在第一对中：有两种位置、\(2m-1\) 种配对指标，以及余下 \((2m-2)!\) 种排列，总计 \(2(2m-1)!\) 项。每项均等于
	\(- (1-t^2)^{m-1}\theta_2\wedge\cdots\wedge\theta_{2m}\)：
	若 \(1\) 在第一位，负号来自 \(A_{1a}=-\theta_a\)；若它在第二位，置换符号再反转一次。其余排列的符号与把楔积重排为 \(2,\ldots,2m\) 的符号相消。因此
	\[
		\operatorname{Pf}(A,\Omega^t,\ldots,\Omega^t)
		=-\frac{(2m-1)!}{2^{m-1}m!}(1-t^2)^{m-1}
		\theta_2\wedge\cdots\wedge\theta_{2m}.
	\]
	当 \(m=1\) 时，曲率因子为空积，计数公式仍给出 \(-\theta_2\)。

	设 \(I_m=\int_0^1(1-t^2)^{m-1}\,\mathrm dt\)。有 \(I_1=1\)，对 \(m\ge2\) 将 \(t(1-t^2)^{m-1}\) 的导数积分，得到
	\[
		(2m-1)I_m=2(m-1)I_{m-1},\qquad
		I_m=\frac{2^{2m-2}((m-1)!)^2}{(2m-1)!}.
	\]
	代入式\eqref{g3:eq:sphere-euler-transgression}，其负号与刚才的负号相消，得到
	\[
		\Psi_\nabla
		=\frac{m}{(2\pi)^m}\,
		\frac{(2m-1)!}{2^{m-1}m!}\,I_m\,
		\theta_2\wedge\cdots\wedge\theta_{2m}
		=\frac{(m-1)!}{2\pi^m}
		\theta_2\wedge\cdots\wedge\theta_{2m}.
	\]
	在 \(u=e_1\) 处，外法向是 \(e_1\)，因此 \((e_2,\ldots,e_{2m})\) 正是球面的正向基，\(\sigma_{2m-1}=\theta_2\wedge\cdots\wedge\theta_{2m}\)。旋转不变性遂给出整根纤维上的式\eqref{g3:eq:sphere-angular-normalization}。

	最后直接核算球面积。逐坐标 Gaussian 积分和极坐标分别给出
	\[
		\pi^m=\int_{\mathbb R^{2m}}e^{-|z|^2}\,\mathrm dz
		=\operatorname{Vol}(S^{2m-1})
		\int_0^\infty e^{-r^2}r^{2m-1}\,\mathrm dr.
	\]
	令 \(s=r^2\)，再分部积分 \(m-1\) 次，径向积分为 \((m-1)!/2\)。所以球面积为 \(2\pi^m/(m-1)!\)，与前面的系数相乘恰为一。
\end{proof}

这说明 \(\Psi_\nabla\) 的纤维限制就是总积分为一的角形式。它在整个球丛上一般不闭，外微分恰好记录底丛的 Euler 形式。若底丛有全局单位截面 \(s:B\to S(E)\)，则拉回后得到
\[
	e(\nabla)=-\mathrm d(s^*\Psi_\nabla).
\]
因而非零截面使 Euler 形式恰当；零点存在时，缺少的原形式将由零点周围的边界积分补足。

\subsection{孤立零点附近的边界积分}
\label{g3:subsec:1702:04}
下面把底空间与丛的秩都取为 \(n=2m\)。设 \(M^n\) 定向，\(E\to M\) 也定向，并有光滑截面 \(s\)，它在点 \(p\) 有孤立零点。在以 \(p\) 为中心的正向坐标球 \(B_\varepsilon\) 上选择正向正交丛标架，写截面系数为 \(F(x)\in\mathbb R^n\)。在 \(\partial B_\varepsilon\) 上归一化，得到
\[
	f_\varepsilon:S^{n-1}\longrightarrow S^{n-1},
	\qquad f_\varepsilon(\theta)
	=\frac{F(\varepsilon\theta)}{|F(\varepsilon\theta)|}.
\]
其映射度称为截面的\term[jiemian jubu zhibiao]{局部指标}[local index]，记为 \(\operatorname{ind}_p(s)\)。源球面和目标球面均取外法向优先方向。改变小半径给出球壳中的归一化同伦；改变正向丛标架时，变换矩阵在球上可沿径向缩到中心的常值正交矩阵，后者的球面度为 \(+1\)。改变正向底坐标，则由第二卷第16.3.1小节中局部度的坐标不变性保持同一个整数。因此指标与这些选择无关。

\ProseParagraphBreak

当 \(E=\mathrm TM\) 并使用与 \(M\) 一致的方向时，这就是第二卷第16.3.1小节的向量场局部指标。那里使用坐标切标架；将它正交化得到正向标架，两个分量映射相差一个处处正行列式的矩阵函数。将矩阵函数缩到中心，再用正行列式线性映射的度为 \(+1\)，可知两种定义相同。第二卷第16.1.2小节的映射度比较定理还给出：对光滑 \(f:S^{n-1}\to S^{n-1}\)，
\begin{equation}\label{g3:eq:angular-degree-integral}
	\int_{S^{n-1}}f^*\left(\frac{\sigma_{n-1}}{\operatorname{Vol}(S^{n-1})}\right)
	=\deg f.
\end{equation}
这里 \(n-1\ge1\)，所以正好适用该处连通闭定向流形的积分公式。

\begin{proposition}\label{g3:prop:angular-local-index-integral}
	设 \(M^{2m}\) 为定向光滑无边界流形，\(m\ge1\)，\(E\to M\) 为定向实秩 \(2m\) Euclid 丛，\(\nabla\) 为相容联络。设光滑截面 \(s\) 在闭坐标球 \(\overline B_\varepsilon\) 内恰有一个零点 \(p\)，且该零点在球内。令 \(\widehat s=s/|s|\)，并用上述构造定义 \(\Psi_\nabla\)。给 \(\partial B_\varepsilon\) 取球的外法向优先方向，则
	\begin{equation}\label{g3:eq:angular-local-index-integral}
		\int_{\partial B_\varepsilon}\widehat s^*\Psi_\nabla
		=\operatorname{ind}_p(s)-\int_{B_\varepsilon}e(\nabla).
	\end{equation}
	特别地，若这些坐标球以 \(p\) 为中心、半径趋零，左边趋于 \(\operatorname{ind}_p(s)\)。零点不要求非退化。
\end{proposition}
\begin{proof}
	先把坐标球写成以 \(p\) 为零点的 \(\{|x|\le\varepsilon\}\)，并在其中取正向正交平凡化。考虑从边界上的单位截面到中心纤维的同伦，写成从中心向边界的参数方向：
	\[
		H:[0,1]\times S^{2m-1}\longrightarrow S(E),\qquad
		H(t,\theta)=(t\varepsilon\theta,f_\varepsilon(\theta)).
	\]
	第二分量表示所选正交标架中的单位向量。即使 \(s\) 在零点非常退化，这个同伦也光滑，因为对固定的 \(\varepsilon\)，只使用非零的边界值。它在 \(t=1\) 等于 \(\widehat s\) 的边界参数化，在 \(t=0\) 则是目标纤维中的 \(f_\varepsilon\)。

	给圆柱取 \(\mathrm dt\) 在前、球面方向在后的乘积取向。其边界在 \(t=1\) 取正号、在 \(t=0\) 取负号，故 Stokes 定理和式\eqref{g3:eq:sphere-euler-transgression-differential}给出
	\[
		\int_{\partial B_\varepsilon}\widehat s^*\Psi_\nabla
		-\int_{S^{2m-1}}f_\varepsilon^*
		\left(\left.\Psi_\nabla\right|_{S(E_p)}\right)
		=-\int_{[0,1]\times S^{2m-1}}(\pi\circ H)^*e(\nabla).
	\]
	第二项由纤维归一化及式\eqref{g3:eq:angular-degree-integral}等于 \(\operatorname{ind}_p(s)\)。底映射 \(\pi\circ H(t,\theta)=t\varepsilon\theta\) 是通常的径向极坐标参数化。\(t>0\) 时，它按 \(\mathrm dt\) 在前的方向正向覆盖穿孔球一次；\(t=0\) 的退化集合不影响顶次积分。因此右边恰为 \(-\int_{B_\varepsilon}e(\nabla)\)，得到所述等式。

	若原球的中心未取在 \(p\)，可用保持边界方向的球内坐标调整将 \(p\) 移到中心；或者先取以 \(p\) 为中心的更小球，再在两球之间对 \(\widehat s^*\Psi_\nabla\) 应用 Stokes 定理，得到相同公式。最后，Euler 形式光滑，在固定坐标中系数有界，因此 \(\int_{B_\varepsilon}e(\nabla)=O(\varepsilon^{2m})\)，极限成立。
\end{proof}

证明中的同伦只移动底点，不改变记录下来的边界方向。这使误差完全由光滑的 Euler 形式在小球内的积分控制，无须对 \(\widehat s\) 的高阶导数作任何估计。因而具有较高绕转数的退化零点，与非退化零点一样适用。

\subsection{Poincaré–Hopf 与整体 Euler 公式}
\label{g3:subsec:1702:05}
现在令 \(M\) 闭。有限个零点给出有限个孔洞，把前一命题的边界贡献相加，就能恢复整个 Euler 积分。

\begin{theorem}\label{g3:thm:euler-form-index-sum}
	设 \(M^{2m}\) 为闭定向光滑流形，\(m\ge1\)，\(E\to M\) 为定向实秩 \(2m\) Euclid 丛，\(\nabla\) 为相容联络。若光滑截面 \(s\) 只有孤立零点，则零点有限，并且
	\begin{equation}\label{g3:eq:euler-form-index-sum}
		\int_Me(\nabla)=\sum_{s(p)=0}\operatorname{ind}_p(s).
	\end{equation}
	允许 \(s\) 处处非零，此时右边为空和。
\end{theorem}
\begin{proof}
	零点集是紧空间中的闭离散集，故有限，记为 \(p_1,\ldots,p_k\)。在各点周围取两两不交的正向坐标球 \(B_i\)，闭球内无其他零点，令
	\[
		M_0=M\setminus\bigcup_{i=1}^k\operatorname{Int}B_i.
	\]
	这里 \(M_0\) 保留了各小球边界，是紧带边界流形。在其上 \(\widehat s\) 全局有定义，且
	\(\mathrm d(\widehat s^*\Psi_\nabla)=-e(\nabla)\)。

	必须区分同一个小球面的两个方向。作为 \(\partial B_i\)，外法向从球内指向外面；作为 \(\partial M_0\)，外法向从保留部分指向孔洞，如\cref{g3:fig:puncture-orientation}所示。这两条法向相反，故
	\[
		\partial M_0=-\coprod_{i=1}^k\partial B_i.
	\]
	因此 Stokes 定理给出
	\[
		-\int_{M_0}e(\nabla)
		=\int_{\partial M_0}\widehat s^*\Psi_\nabla
		=-\sum_{i=1}^k\int_{\partial B_i}\widehat s^*\Psi_\nabla.
	\]
	代入式\eqref{g3:eq:angular-local-index-integral}并移项，得到
	\[
		\int_{M_0}e(\nabla)+\sum_i\int_{B_i}e(\nabla)
		=\sum_i\operatorname{ind}_{p_i}(s).
	\]
	左边就是 \(\int_Me(\nabla)\)，所以甚至无须再令球半径趋零。若没有零点，直接在闭流形上积分恰当形式
	\(e(\nabla)=-\mathrm d(\widehat s^*\Psi_\nabla)\)，也得到零。
\end{proof}

\begin{figure}[htbp]
	\centering
	% !TeX root = ../Book3.tex
% 两幅图只画一个坐标邻域；灰色区域没有画出人为外边界。
\begin{tikzpicture}[x=1cm,y=1cm,>=stealth,font=\small]
	\begin{scope}[shift={(-2.5,0)}]
		\fill[blue!9] (0,0) circle (1);
		\draw[thick,blue!70!black] (0,0) circle (1);
		\draw[very thick,blue!70!black,->] (25:1) arc[start angle=25,end angle=88,radius=1];
		\draw[thick,->] (1,0)--(1.7,0) node[right] {\(\nu_B\)};
		\fill (0,0) circle (1.2pt) node[below] {\(p\)};
		\node at (0,.43) {\(B\)};
		\node[below] at (0,-1.25) {\(\partial B\)：球的边界方向};
	\end{scope}
	\begin{scope}[shift={(2.7,0)}]
		\fill[gray!13] (-1.65,-1.15) rectangle (1.65,1.15);
		\fill[white] (0,0) circle (1);
		\draw[thick,red!75!black] (0,0) circle (1);
		\draw[very thick,red!75!black,->] (145:1) arc[start angle=145,end angle=82,radius=1];
		\draw[thick,->] (1,0)--(.37,0) node[below left] {\(\nu_{M_0}\)};
		\fill (0,0) circle (1.2pt) node[below] {\(p\)};
		\node at (0,.43) {挖去的球};
		\node[right] at (1.1,.8) {\(M_0\)};
		\node[below] at (0,-1.25) {\(\partial M_0\)：朝孔洞的外法向};
	\end{scope}
\end{tikzpicture}

	\caption{同一小球面的两种边界方向。右图只画孔洞附近的坐标区域，灰色表示保留部分，并非给闭流形添加一个外边界。}
	\label{g3:fig:puncture-orientation}
\end{figure}

切丛的局部指标和已由第二卷第17.4.3小节的 Poincaré--Hopf 定理确定：在任意正维光滑闭流形上，只有孤立零点的光滑向量场满足
\[
	\sum_{X(p)=0}\operatorname{ind}_p(X)=\chi(M).
\]
该处先用 Morse 梯度场的局部指标 \((-1)^\lambda\) 和把手的 Euler 等式得到结论，再用总指标的同伦不变性推广到任意孤立零点场。这里需要的正是闭流形版本。第二卷第17.1.3小节还证明了光滑 Morse 函数的存在；在闭流形上，其梯度场只有有限个非退化零点，因而总能提供满足要求的向量场。

\begin{theorem}[Gauss--Bonnet--陈定理]\label{g3:thm:gauss-bonnet-chern}
	设 \((M^{2m},g)\) 为闭定向光滑黎曼流形，\(m\ge1\)，\(\nabla^g\) 为 Levi--Civita 联络。在正向正交标架中写
	\(\nabla^ge_j=\sum_i e_i\omega_{ij}\)，并记
	\(\Omega=\mathrm d\omega+\omega\wedge\omega\)。则
	\begin{equation}\label{g3:eq:gauss-bonnet-chern}
		\frac1{(2\pi)^m}\int_M\operatorname{Pf}(\Omega)=\chi(M).
	\end{equation}
	同一公式也适用于 \(\mathrm TM\) 上任意保持 \(g\) 的联络，不要求该联络无挠。
\end{theorem}
\begin{proof}
	选取上述 Morse 梯度场 \(X\)，将\cref{g3:thm:euler-form-index-sum}用于 \(E=\mathrm TM\) 及给定的相容联络。切丛的方向与 \(M\) 的方向一致，所以局部指标正是 Poincaré--Hopf 定理所用的指标。于是
	\[
		\frac1{(2\pi)^m}\int_M\operatorname{Pf}(\Omega)
		=\int_Me(\nabla)
		=\sum_{X(p)=0}\operatorname{ind}_p(X)
		=\chi(M).
	\]
	证明只用了度量相容性，没有使用无挠性，因此最后一条断言同时成立。不连通时，在有限个连通分支上分别应用并相加即可。
\end{proof}

式\eqref{g3:eq:gauss-bonnet-chern}称为\term[Gauss Bonnet Chen dingli]{Gauss--Bonnet--陈定理}[Gauss--Bonnet--Chern theorem]。其整数值来自球面映射度，纤维积分的归一化则把这个整数与具体曲率多项式连接起来。若反转 \(M\) 的方向，切丛的方向及 Pfaffian 同时反转，积分所用方向也反转一次，所以最终数值不变。

\begin{example}\label{g3:exm:euler-surface-and-flat}
	二维时 \(\Omega_{12}=K\,\mathrm dA\)，所以公式恢复
	\[
		\frac1{2\pi}\int_MK\,\mathrm dA=\chi(M).
	\]
	在任意偶维，若闭定向黎曼流形的 Levi--Civita 曲率恒为零，Pfaffian 也恒为零，故其 Euler 示性数为零。更一般地，只要定向切丛允许某个平坦的度量联络，同一结论就成立；此处的联络可以不同于 Levi--Civita 联络。
\end{example}

闭性在 Stokes 步骤中有明确作用。对于带边界流形，原有外边界不会与零点小球相消，会留下单位向量场沿边界的转移积分；习题将推导这个版本。不可定向时，局部 Euler 形式要与局部体积方向同时处理，才能拼成可积分密度；其定向双覆盖版本也在习题中给出。
% 纲要条目排为编号小节标题，并进入目录与书签；说明段仍保留于纲要注释中。

% !TeX root = ../../Book3.tex
% 纲要标题：### 17.3* Dirac 型算子与自旋
% 纲要位置：几何学计划纲要.md，第 2800 行。
% 正文状态：初稿；尚未正式审核。

% 纲要内容开始
% 
% * Clifford 代数和 Clifford 丛
% * spin 提升与自旋子丛
% * Dirac 算子及符号
% * Lichnerowicz 公式与正标量曲率障碍
% 
% 纲要内容结束

\OptionalSection{Dirac 型算子与自旋}{g3:sec:1703}

Laplace 算子的主符号是 \(|\xi|^2\) 乘以恒等映射。如果希望找到一个一阶算子，使它的平方具有这个主符号，就需要把切向量表示成满足特定反交换关系的线性映射。Clifford 代数记录这一关系，自旋结构则解决这些局部表示能否沿流形一致粘合的问题。本节把代数构造、联络和椭圆算子接在一起，最后从一个点态恒等式推出正标量曲率的障碍。

\ProseParagraphBreak

本节的黎曼流形均光滑、无边界，维数 \(n\ge2\)；需要紧性时另行说明。曲率沿用
\(R(X,Y)=[\nabla_X,\nabla_Y]-\nabla_{[X,Y]}\)，Clifford 乘法采用负号约定，Hermitian 内积对第一变量线性。第二卷已建立 spin 提升的拓扑判据；这里还要具体构造提升后的表示和微分算子。

\subsection{Clifford 代数和 Clifford 丛}
\label{g3:subsec:1703:01}
设 \((V,g)\) 是 \(n\) 维实内积空间。其\term[Clifford daishu]{Clifford 代数}[Clifford algebra]定义为张量代数的商
\[
	\operatorname{Cl}(V,g)=T(V)\big/\langle v\otimes v+g(v,v)1:v\in V\rangle.
\]
由极化，商中的基本关系为
\begin{equation}\label{g3:eq:clifford-negative-relation}
	vw+wv=-2g(v,w)1.
\end{equation}
因此，给定含幺代数 \(A\) 以及满足同一关系的线性映射 \(c:V\to A\)，就存在唯一含幺代数同态 \(\operatorname{Cl}(V,g)\to A\) 延拓 \(c\)。这种把二次关系放入乘法的办法源自 Clifford 的工作。\footnote{William Kingdon Clifford（1845--1879），英国数学家，研究非欧几何以及几何与代数的联系，提出了后来以他命名的代数。}

\begin{proposition}\label{g3:prop:clifford-basis}
	设 \(e_1,\ldots,e_n\) 是实内积空间 \(V\) 的正交归一基。记 \(e_\varnothing=1\)，并对 \(I=\{i_1<\cdots<i_r\}\) 记 \(e_I=e_{i_1}\cdots e_{i_r}\)。则这些 \(2^n\) 个元素构成 \(\operatorname{Cl}(V,g)\) 的一组实基。奇、偶长度单项式分别张成其 \(\mathbb Z/2\mathbb Z\) 分次的奇、偶部分。
\end{proposition}
\begin{proof}
	关系 \(e_i^2=-1\) 和 \(e_ie_j=-e_je_i\)（\(i\ne j\)）把每个单项式整理成某个 \(e_I\) 的带符号倍数，所以这些元素张成商代数。
	为证线性无关，在 \(\Lambda^*V^*\) 上令
	\[
		c(v)=\varepsilon(v^\flat)-\iota_v,
	\]
	其中 \(\varepsilon\) 表示外乘。恒等式
	\(\iota_v\varepsilon(w^\flat)+\varepsilon(w^\flat)\iota_v=g(v,w)\)
	和楔乘、收缩各自的反交换关系给出
	\(c(v)c(w)+c(w)c(v)=-2g(v,w)\)。于是得到商代数的表示。将 \(c(e_{i_1})\cdots c(e_{i_r})\) 作用于常数 \(1\)，由于各指标不同，收缩项全为零，结果是 \(e^{i_1}\wedge\cdots\wedge e^{i_r}\)。这些外形式线性无关，故 \(e_I\) 也线性无关。定义关系只改变偶数个因子，因而乘积保留长度的奇偶性，给出所述分次。
\end{proof}

在实 Clifford 代数上有反自同构 \(a\mapsto a^*\)，由 \(v^*=-v\) 和 \((ab)^*=b^*a^*\) 确定；在复化 \(\operatorname{Cl}^{\mathbb C}(V)=\operatorname{Cl}(V)\otimes_{\mathbb R}\mathbb C\) 上再令它对标量共轭。正交映射保留定义关系，故自然作用于 Clifford 代数。把这个作用用于切丛的正交标架转移函数，便得到光滑代数丛 \(\operatorname{Cl}(TM)\)。其存在不要求流形定向，也不要求自旋结构。Levi--Civita 联络按乘积法则延拓到此丛；由于 \(\nabla g=0\)，它保留商关系。

\ProseParagraphBreak

自旋子所用的表示比 \(\Lambda^*V^*\) 小。在偶数维 \(n=2m\) 中，取
\(\Sigma_{2m}=\Lambda^*\mathbb C^m\)，并令 \(\varepsilon_a\) 为第 \(a\) 个标准基向量的外乘，\(\iota_a=\varepsilon_a^*\) 为对应收缩。定义
\begin{equation}\label{g3:eq:spin-module-explicit}
	c(e_{2a-1})=\varepsilon_a-\iota_a,
	\qquad c(e_{2a})=\mathrm i(\varepsilon_a+\iota_a),
	\qquad 1\le a\le m.
\end{equation}
每个 \(c(e_j)\) 都是斜 Hermitian 映射；直接使用楔乘与收缩关系可得 \(c(e_j)^2=-1\)，不同指标的映射反交换。因此它给出 \(\operatorname{Cl}_{2m}^{\mathbb C}\) 的表示。

\begin{proposition}\label{g3:prop:complex-spin-module}
	当 \(n=2m\ge2\) 时，式\eqref{g3:eq:spin-module-explicit}诱导代数同构
	\[
		\operatorname{Cl}_{2m}^{\mathbb C}\simeq\operatorname{End}_{\mathbb C}(\Sigma_{2m}),
		\qquad \dim_{\mathbb C}\Sigma_{2m}=2^m.
	\]
	定义
	\begin{equation}\label{g3:eq:spin-chirality}
		\Gamma=\mathrm i^m c(e_1)\cdots c(e_{2m}).
	\end{equation}
	则 \(\Gamma^*=\Gamma\)、\(\Gamma^2=1\)，它在 \(\Lambda^r\mathbb C^m\) 上等于 \((-1)^r\)。它与每个 \(c(v)\) 反交换，与偶 Clifford 元素的作用交换。两特征空间 \(\Sigma_{2m}^{\pm}\) 的维数均为 \(2^{m-1}\)。
\end{proposition}
\begin{proof}
	由式\eqref{g3:eq:spin-module-explicit}可以线性解出所有 \(\varepsilon_a,\iota_a\)。算子
	\(P_0=\prod_{a=1}^m\iota_a\varepsilon_a\)
	是到 \(\Lambda^0\mathbb C^m\) 的投影。若 \(I=\{i_1<\cdots<i_r\}\)、\(J=\{j_1<\cdots<j_s\}\)，算子
	\(\varepsilon_{i_1}\cdots\varepsilon_{i_r}P_0\iota_{j_s}\cdots\iota_{j_1}\)
	把基向量 \(e_J\) 送到 \(e_I\)，把其余外代数基向量送到零。因此表示的像含有全部矩阵单位，即全部端同态代数。两边复维数均为 \(2^{2m}\)，故它是同构，特别地这个表示不可约。
	记 \(N_a=\varepsilon_a\iota_a\)，则
	\(\mathrm i c(e_{2a-1})c(e_{2a})=1-2N_a\)。不同 \(a\) 的这些偶算子交换，故 \(\Gamma=\prod_a(1-2N_a)\)，这正是外次数的奇偶算子。其余性质随即成立；\(c(e_1)\) 在两个特征空间间给出同构，所以二者维数相等。
\end{proof}

若 \(n=2m+1\ge3\)，仍取 \(\Sigma_n=\Lambda^*\mathbb C^m\)，保留前 \(2m\) 个作用，再定义 \(c(e_{2m+1})=\mathrm i\Gamma\)。这个新作用平方为 \(-1\)，并与其余作用反交换，故同样得到不可约复 Clifford 模，维数为 \(2^m\)。也可选择相反号，以下固定上述选择。奇数维中，前 \(2m\) 个方向产生的 \(\Gamma\) 依赖这个方向的选取；它不是流形上的自然手性分次。本节只在偶数维中使用 \(S^+\oplus S^-\)。

\subsection{spin 提升与自旋子丛}
\label{g3:subsec:1703:02}
记 \(\operatorname{Cl}_n=\operatorname{Cl}(\mathbb R^n)\)。定义 \(\operatorname{Spin}(n)\) 为其中所有偶数个单位向量的乘积所成的群，空乘积等于 \(1\)。单位向量 \(u\) 的逆为 \(-u\)，所以这确实是群，且包含于偶部分。它通过共轭作用于 \(\mathbb R^n\)：
\[
	\rho(s)v=svs^{-1}.
\]
下面同时核对这一作用的值域和它的覆盖性质。

\begin{theorem}\label{g3:thm:spin-double-cover}
	对每个 \(n\ge2\)，\(\operatorname{Spin}(n)\) 是连通 Lie 群，\(\rho:\operatorname{Spin}(n)\to\operatorname{SO}(n)\) 是满的二重覆盖，其核为 \(\{1,-1\}\)。当 \(n=2\) 时，
	\[
		\operatorname{Spin}(2)=\{\cos t+\sin t\,e_1e_2:t\in\mathbb R\},
	\]
	且 \(\rho\) 把参数 \(t\) 送到平面上角度 \(2t\) 的旋转。
\end{theorem}
\begin{proof}
	设 \(r_u(v)=v-2g(u,v)u\) 是单位向量 \(u\) 的正交超平面反射。Clifford 关系给出 \(uvu^{-1}=-r_u(v)\)。于是偶数个单位向量的共轭作用是偶数个反射的乘积，属于 \(\operatorname{SO}(n)\)。每个正交映射都是反射的乘积：若 \(Ae_1\ne e_1\)，以 \(Ae_1-e_1\) 为法向的反射将 \(Ae_1\) 送到 \(e_1\)；在 \(e_1^\perp\) 上继续归纳，至多用 \(n\) 次反射就把 \(A\) 化为恒等映射。若 \(Ae_1=e_1\)，直接在正交补上归纳。行列式等于 \(1\) 时反射的个数必为偶数。因此 \(\rho\) 满。
	若偶元素 \(s\) 与所有 \(e_j\) 交换，把它写成有序单项式的线性组合。共轭 \(e_j(\cdot)e_j^{-1}\) 在偶单项式 \(e_I\) 上的作用，当 \(j\notin I\) 时为 \(+1\)，当 \(j\in I\) 时为 \(-1\)。所以同时与全部 \(e_j\) 交换的偶元素只能是标量。又因单位向量的乘积满足 \(s^*s=1\)，标量只能为 \(\pm1\)。两者确实都在群中，因为 \(e_1e_1=-1\)。这证明核的断言。
	还需证明它在单位元附近确实是光滑覆盖。对斜对称映射 \(A:\mathbb R^n\to\mathbb R^n\)，令
	\begin{equation}\label{g3:eq:spin-lie-algebra-lift}
		b(A)=\frac14\sum_{i,j=1}^n g(Ae_i,e_j)e_ie_j.
	\end{equation}
	从
	\([e_ie_j,e_k]=2(\delta_{ik}e_j-\delta_{jk}e_i)\)
	可得 \([b(A),v]=Av\)。在正交二维旋转块中，\(b(A)\) 是若干彼此交换的 \((\theta/2)e_ie_j\) 之和；而
	\[
		\exp(te_ie_j)=\cos t+\sin t\,e_ie_j
		=e_i(-\cos t\,e_i+\sin t\,e_j)
	\]
	是两个单位向量的乘积。因此 \(\exp b(A)\in\operatorname{Spin}(n)\)。对共轭作用求导并使用 \([b(A),v]=Av\)，得到 \(\rho(\exp b(A))=\exp A\)。在 \(\operatorname{SO}(n)\) 单位元的充分小邻域 \(U\) 内，矩阵对数光滑，故
	\[
		h(Q)=\exp b(\log Q),\qquad Q\in U,
	\]
	是光滑截面。核的计算说明 \(\rho^{-1}(U)=h(U)\sqcup(-h(U))\)。缩小 \(U\)，两个分支分别位于 Clifford 代数中 \(1\) 与 \(-1\) 的不交邻域；它们在 \(\rho^{-1}(U)\) 中都是开集，且 \(\rho\) 分别与 \(h\)、\(-h\) 互逆。这些局部图经群平移给出 Lie 群结构和二重覆盖；乘法与求逆在这些图中光滑，因为它们来自有限维 Clifford 代数的乘法及 \(s^{-1}=s^*\)。
	最后，\(S^{n-1}\) 在 \(n\ge2\) 时连通。任一单位向量乘积 \(uv\) 都可沿 \(v\) 到 \(-u\) 的路径连接到 \(u(-u)=1\)；把一个偶乘积分成这样的二因子组，即知群连通。二维时偶代数由 \(1,e_1e_2\) 张成，且 \((e_1e_2)^2=-1\)，上面的指数公式给出整个群。共轭在 \(e_1\) 上的值为 \(\cos(2t)e_1+\sin(2t)e_2\)，故旋转角确为 \(2t\)。
\end{proof}

二维的覆盖是圆周的平方映射，并非普适覆盖；在 \(n\ge3\) 时，第二卷第8.4.3小节的“经典群的一、二次同伦群”命题给出 \(\pi_1(\operatorname{SO}(n))=\mathbb Z/2\mathbb Z\)，所以连通二重覆盖就是普适覆盖。这一区别不影响下面的提升定义。

\begin{definition}\label{g3:def:riemannian-spin-structure}
	设 \((M^n,g)\) 是定向黎曼流形，\(n\ge2\)，\(P_{\operatorname{SO}}(TM)\) 是其定向正交标架主丛。一个\term[spin-jiegou]{spin 结构}[spin structure]（也称自旋结构）由主 \(\operatorname{Spin}(n)\) 丛 \(Q\to M\) 和覆盖恒等映射的光滑映射 \(\eta:Q\to P_{\operatorname{SO}}(TM)\) 组成，满足
	\[
		\eta(qs)=\eta(q)\rho(s),
	\]
	且 \(\eta\) 在每个纤维上是上述二重覆盖。等价结构的丛同构还须与 \(\eta\) 相容。
\end{definition}

第二卷第21.4.4小节的“spin 提升准则”已证明：对紧光滑流形上的定向实向量丛 \(E\)（秩至少为 \(2\)），这样的提升存在当且仅当 \(w_2(E)=0\)；存在时，等价类集合是 \(H^1(M;\mathbb F_2)\) 的挠子。其具体含义也可以从转移函数看出：把 \(\operatorname{SO}(n)\) 的转移函数逐个提升到 \(\operatorname{Spin}(n)\) 后，三重交叠上的乘积可能为 \(-1\)；这些符号构成代表 \(w_2\) 的二上闭链。提升存在后，改变粘合的符号得到一上闭链，等价改变对应一上边界。因此自旋结构是一项额外的整体选择。

\ProseParagraphBreak

例如，第二卷第21.4.2小节“球丛和低维切丛”中的闭定向曲面例已算出：\(w_2(TM)\) 是 Euler 类模 \(2\) 的约化，而 Euler 数 \(2-2g\) 为偶数，所以 \(w_2(TM)=0\)。结合第21.4.4小节的提升及分类准则，亏格 \(g\) 的闭连通定向曲面便有 \(2^{2g}\) 个自旋结构等价类。相反，第21.4.4小节的复射影空间例给出 \(w_2(T\mathbb{CP}^2)\ne0\)，所以 \(\mathbb{CP}^2\) 没有自旋结构。前述 \(\operatorname{Cl}(TM)\) 在这两个情形中却都存在。

\ProseParagraphBreak

把第一小节的表示限制到 \(\operatorname{Spin}(n)\)，记为 \(\kappa\)。由于每个单位向量的 Clifford 作用都是酉映射，\(\kappa\) 是酉表示。给定自旋结构，定义其复\term[zixuanzi cong]{自旋子丛}[spinor bundle]
\[
	S=Q\times_\kappa\Sigma_n,
	\qquad (qs,z)\sim(q,\kappa(s)z).
\]
它的复秩是 \(2^{\lfloor n/2\rfloor}\)。若 \(v\in T_xM\)，\(q\in Q_x\) 对应标架 \(\eta(q):\mathbb R^n\to T_xM\)，令
\[
	c(v)[q,z]=[q,c(\eta(q)^{-1}v)z].
\]
恒等式 \(\kappa(s)c(a)\kappa(s)^{-1}=c(\rho(s)a)\) 保证这一定义与代表无关。偶数维时，定向体积元素及 \(\Gamma\) 在 \(\operatorname{Spin}(n)\) 下不变，因此得到全局正交分解
\(S=S^+\oplus S^-\)，两丛的秩均为 \(2^{n/2-1}\)。

\begin{proposition}\label{g3:prop:spin-connection-curvature}
	设 \((M^n,g)\) 是定向黎曼自旋流形，\(n\ge2\)。Levi--Civita 联络唯一提升为 \(Q\) 上的主联络，并在 \(S\) 上诱导 Hermitian 联络 \(\nabla^S\)，满足
	\begin{equation}\label{g3:eq:spin-connection-compatibility}
		\nabla^S_X(c(v)\psi)=c(\nabla_Xv)\psi+c(v)\nabla^S_X\psi.
	\end{equation}
	在一个可提升的定向正交标架 \(e_1,\ldots,e_n\) 中，令 \(c_a=c(e_a)\)，则其局部公式和曲率公式为
	\begin{align}
		\nabla^S_X\psi
		&=X\psi+\frac14\sum_{a,b}g(\nabla_Xe_a,e_b)c_ac_b\psi,
		\label{g3:eq:spin-connection-local}\\
		R^S(X,Y)
		&=\frac14\sum_{a,b}g(R(X,Y)e_a,e_b)c_ac_b.
		\label{g3:eq:spin-curvature-local}
	\end{align}
	在偶数维中，此联络保留 \(S^+\) 和 \(S^-\)。
\end{proposition}
\begin{proof}
	由二重覆盖可知 \(\rho_*:\mathfrak{spin}(n)\to\mathfrak{so}(n)\) 是同构，其逆恰是式\eqref{g3:eq:spin-lie-algebra-lift}的 \(b\)。若 \(\omega\) 是正交标架丛的 Levi--Civita 主联络形式，则 \(b(\eta^*\omega)\) 满足主联络的等变条件和在竖直向量上的归一化条件，所以给出提升。任何提升都必须经 \(\rho_*\) 送到 \(\eta^*\omega\)，故唯一。代入表示即得式\eqref{g3:eq:spin-connection-local}。
	对实斜对称系数，\(c_ac_b\)（\(a\ne b\)）是斜 Hermitian 映射，所以该联络保留内积。已证的交换子关系 \([b(A),v]=Av\) 在表示中变为 \([c(b(A)),c(v)]=c(Av)\)；加上对 \(v\) 的系数求导，恰好给出式\eqref{g3:eq:spin-connection-compatibility}。偶 Clifford 作用与 \(\Gamma\) 交换，故分次也被保留。
	最后，Lie 群同态的微分保留括号，故提升联络的曲率为 \(b\) 作用于原联络曲率。也可直接把式\eqref{g3:eq:spin-connection-local}代入两个协变导数的交换子：一阶项相消，零阶项就是 \(c(b(R(X,Y)))\)。由 \(b\) 的公式即得式\eqref{g3:eq:spin-curvature-local}。
\end{proof}

按本章的列指标记号 \(\nabla e_b=\sum_a e_a\omega_{ab}\)，有 \(\omega_{ab}(X)=g(\nabla_Xe_b,e_a)\)，所以式\eqref{g3:eq:spin-connection-local}中的联络项也可写为 \(-\tfrac14\sum_{a,b}\omega_{ab}(X)c_ac_b\)。曲率公式在相同列指标约定下具有同样的负号。

\subsection{Dirac 算子及符号}
\label{g3:subsec:1703:03}
设 \(E\to M\) 是 Hermitian 丛，带有满足式\eqref{g3:eq:clifford-negative-relation}的斜 Hermitian 作用 \(c:TM\to\operatorname{End}(E)\)，并带有满足式\eqref{g3:eq:spin-connection-compatibility}同类乘积法则的 Hermitian 联络 \(\nabla^E\)。这种资料称为带相容联络的 Clifford 模丛。其\term[Dirac suanzi]{Dirac 算子}[Dirac operator]定义为
\begin{equation}\label{g3:eq:clifford-dirac-definition}
	D_E\psi=\sum_{j=1}^n c(e_j)\nabla^E_{e_j}\psi.
\end{equation}
这是把协变导数与 Clifford 乘法作收缩，所以不依赖正交标架。取 \(E=S\) 及其自旋联络时，称为自旋 Dirac 算子，以下简记为 \(D\)。它的名称来自 Dirac 对相对论性波动方程的一阶化。\footnote{Paul Adrien Maurice Dirac（1902--1984），英国理论物理学家，参与建立量子力学的数学形式，并于1928年提出描述相对论性自旋 \(1/2\) 粒子的 Dirac 方程。}

\begin{theorem}\label{g3:thm:spin-dirac-elliptic}
	设 \((M^n,g)\) 是光滑无边界黎曼流形，\(n\ge2\)，\(E\) 是上述带相容 Hermitian 联络的 Clifford 模丛。则 \(D_E\) 的主符号及其平方为
	\begin{equation}\label{g3:eq:dirac-principal-symbol}
		\sigma_1(D_E)(x,\xi)=\mathrm i c(\xi^\sharp),
		\qquad \sigma_1(D_E)(x,\xi)^2=|\xi|^2\operatorname{Id}_{E_x}.
	\end{equation}
	因此 \(D_E\) 椭圆且形式自伴。若 \(M\) 闭，则它在 \(L^2(E)\) 上以 \(H^1(E)\) 为定义域是自伴算子，核有限维且由光滑截面组成；其 \(H^1(E)\to L^2(E)\) 实现是 Fredholm 算子。
\end{theorem}
\begin{proof}
	由乘积法则，\(D_E(f\psi)-fD_E\psi=c((\mathrm df)^\sharp)\psi\)。本卷的符号约定把一阶导数替换为 \(\mathrm i\xi\)，故得到第一式。再用 Clifford 关系，\((\mathrm i c(\xi^\sharp))^2=|\xi|^2\)，这同时证明非零余切向量处的可逆性。
	为证形式自伴，取紧支光滑截面 \(\psi,\varphi\)，定义复向量场
	\(Z=\sum_j\langle c(e_j)\psi,\varphi\rangle e_j\)。在任一点选取正规正交标架，利用联络相容性和 \(c(e_j)^*=-c(e_j)\)，得到
	\[
		\operatorname{div}Z
		=\sum_j\bigl(\langle c(e_j)\nabla^E_{e_j}\psi,\varphi\rangle
		+\langle c(e_j)\psi,\nabla^E_{e_j}\varphi\rangle\bigr)
		=\langle D_E\psi,\varphi\rangle-\langle\psi,D_E\varphi\rangle.
	\]
	对实部、虚部分别应用紧支散度公式，即得 \((D_E\psi,\varphi)_{L^2}=(\psi,D_E\varphi)_{L^2}\)。闭流形上的定义域和自伴性由\cref{g3:prop:elliptic-adjoint-domain}推出；有限维光滑核与 Fredholm 性由\cref{g3:thm:elliptic-fredholm}推出。
\end{proof}

自旋 Dirac 算子的平方是非负的二阶算子：闭流形上
\((D^2\psi,\psi)=\|D\psi\|_2^2\)。在偶数维 \(n=2m\) 中，\(\Gamma\) 平行且与 Clifford 向量作用反交换，所以
\[
	D\Gamma=-\Gamma D,\qquad
	D=\begin{pmatrix}0&D^-\\D^+&0\end{pmatrix}
	\quad\text{于 }S^+\oplus S^-.
\]
\(\mathrm i c(\xi^\sharp)\) 在 \(S^+\) 与 \(S^-\) 之间互为同构，故 \(D^+:H^1(S^+)\to L^2(S^-)\) 也椭圆。其形式伴随为 \(D^-\)，因此在闭流形上
\begin{equation}\label{g3:eq:spin-analytic-index}
	\operatorname{ind}D^+=\dim\ker D^+-\dim\ker D^-.
\end{equation}
这里全算子 \(D\) 的自伴性与分次算子 \(D^+\) 的指标是两个相容的事实：前者的指标总为零，后者比较的是两个不同手性中的零模数。

\ProseParagraphBreak

在欧氏空间及其平坦自旋标架中，\(D=\sum_jc(e_j)\partial_j\)，反交换关系立即给出 \(D^2=-\sum_j\partial_j^2\)。曲率非零时，协变导数彼此不交换。下一小节将计算这个交换子产生的全部零阶项。

\subsection{Lichnerowicz 公式与正标量曲率障碍}
\label{g3:subsec:1703:04}
先在任意带相容联络的 Clifford 模丛上计算。于一点选取满足 \(\nabla e_j=0\) 的正规正交标架，展开式\eqref{g3:eq:clifford-dirac-definition}并把不同指标配对，得到
\[
	D_E^2=-\sum_i\nabla^E_{e_i}\nabla^E_{e_i}
	+\sum_{i<j}c(e_i)c(e_j)R^E(e_i,e_j).
\]
恢复粗 Laplace 算子的坐标不变表达式，即为
\begin{equation}\label{g3:eq:dirac-square-general-curvature}
	D_E^2=(\nabla^E)^*\nabla^E
	+\frac12\sum_{i,j}c(e_i)c(e_j)R^E(e_i,e_j).
\end{equation}
联络曲率在这一步尚未化成标量。对自旋子丛，式\eqref{g3:eq:spin-curvature-local}使这一收缩特别简单。

\begin{theorem}[Lichnerowicz 公式]\label{g3:thm:spin-lichnerowicz}
	设 \((M^n,g)\) 是光滑无边界定向黎曼自旋流形，\(n\ge2\)，\(S\) 配备 Levi--Civita 联络诱导的自旋联络。则其自旋 Dirac 算子满足
	\begin{equation}\label{g3:eq:spin-lichnerowicz}
		D^2=(\nabla^S)^*\nabla^S+\frac{\operatorname{Scal}_g}{4}\operatorname{Id}_S.
	\end{equation}
	这是局部微分算子恒等式，不要求 \(M\) 紧。
\end{theorem}
\begin{proof}
	为把本卷的曲率符号和 Clifford 乘法同时保留，记
	\[
		C_{ijab}=g(R(e_i,e_j)e_a,e_b),\qquad c_j=c(e_j).
	\]
	它在 \((i,j)\)、\((a,b)\) 内分别反对称，且 \(C_{ijab}=C_{abij}\)。注意这里最后两个指标的顺序与第三章的 \(\operatorname{Rm}(e_i,e_j,e_a,e_b)\) 相反。由式\eqref{g3:eq:dirac-square-general-curvature}和式\eqref{g3:eq:spin-curvature-local}，待计算的零阶项为
	\[
		\mathcal K=\frac18\sum_{i,j,a,b}C_{ijab}c_ic_jc_ac_b.
	\]
	在每一对指标相同的项本来就是零，因此非零项的指标集只有两、三或四个不同元素。下面分别归并。
	若有四个不同指标，固定 \(i<j<a<b\)，把它们按次序移到 \(c_ic_jc_ac_b\) 后，系数为
	\[
		8\bigl(C_{ijab}-C_{iajb}+C_{ibja}\bigr).
	\]
	因子 \(8\) 来自每一对内部的两种次序和两对的交换。第一 Bianchi 恒等式给出
	\(C_{ijab}+C_{jaib}+C_{aijb}=0\)；再用对称性将后两项分别写成 \(C_{ibja}\) 与 \(-C_{iajb}\)，故这一系数为零。
	若恰有三个不同指标，两对指标共有一个。交换两对不改变 \(C_{ijab}\)，而相应的两个 Clifford 二重积反交换。例如在三个指标互异时，
	\[
		(c_ic_j)(c_ic_b)+(c_ic_b)(c_ic_j)=0.
	\]
	其余共有指标的位置通过两对内的反对称性化到这一式。因此这些项也成对抵消。
	最后，若只有两个不同指标，则对每个有序对 \(i\ne j\)，只有 \((a,b)=(i,j)\) 或 \((j,i)\) 两项。由 \(c_ic_jc_ic_j=-1\) 和 \(c_ic_jc_jc_i=1\)，它们的和为
	\[
		-2C_{ijij}\operatorname{Id}_S
		=2g(R(e_i,e_j)e_j,e_i)\operatorname{Id}_S.
	\]
	而 \(\sum_{i,j}g(R(e_i,e_j)e_j,e_i)=\operatorname{Scal}_g\)。故总和为 \(2\operatorname{Scal}_g\operatorname{Id}_S\)，再乘 \(1/8\) 即得 \(\mathcal K=\operatorname{Scal}_g/4\)，证明公式。
\end{proof}

这个计算解释了为什么自旋子最终只看到标量曲率：四个不同指标由 Bianchi 恒等式消去，三个不同指标由 Clifford 反交换消去，留下的两指标收缩正是截面曲率之和。一般 Clifford 模丛仍应使用式\eqref{g3:eq:dirac-square-general-curvature}；其联络曲率未必具有自旋曲率公式。

\begin{corollary}\label{g3:cor:spin-positive-scalar-vanishing}
	设 \((M^n,g)\) 是非空闭连通定向黎曼自旋流形，\(n\ge2\)。若 \(\operatorname{Scal}_g\ge0\)，则每个满足 \(D\psi=0\) 的自旋子都是平行的。若标量曲率还在某一点严格为正，则 \(\ker D=0\)。特别地，若 \(\operatorname{Scal}_g\ge s_0>0\)，则对一切光滑自旋子有
	\[
		\|D\psi\|_2^2\ge\frac{s_0}{4}\|\psi\|_2^2.
	\]
	在偶数维中，只要 \(g\) 的标量曲率处处为正，就有 \(\operatorname{ind}D^+=0\)。
\end{corollary}
\begin{proof}
	积分式\eqref{g3:eq:spin-lichnerowicz}并分部积分，得到
	\begin{equation}\label{g3:eq:spin-integrated-lichnerowicz}
		\|D\psi\|_2^2=\|\nabla^S\psi\|_2^2
		+\frac14\int_M\operatorname{Scal}_g|\psi|^2\,\mu_g.
	\end{equation}
	当 \(D\psi=0\) 且标量曲率非负时，两项都为零，故 \(\psi\) 平行。平行截面的长度在连通流形上为常数；如果标量曲率在一点为正，其与 \(|\psi|^2\) 的乘积在那里为零，就迫使这个常数为零。下界估计也直接来自积分公式。偶数维时 \(\ker D=\ker D^+\oplus\ker D^-\)，所以两个核均为零，式\eqref{g3:eq:spin-analytic-index}给出指标为零。
\end{proof}

为了把“给定度量的核消失”变成“流形不允许某种度量”，还要证明所用的指标不会随度量改变。这里需要识别不同度量的自旋子丛，而不能直接把它们当成同一个丛。

\begin{proposition}\label{g3:prop:spin-index-metric-independent}
	设 \(M^{2m}\) 是闭定向光滑流形，\(m\ge1\)，固定其一个自旋结构。按下面的标架识别沿度量变化运输此结构，则 \(\operatorname{ind}D_g^+\) 与黎曼度量 \(g\) 无关。因此，如果这个指标非零，\(M\) 不存在处处正标量曲率的黎曼度量。
\end{proposition}
\begin{proof}
	给定两度量 \(g_0,g_1\)，取 \(g_t=(1-t)g_0+tg_1\)。定义正定 \(g_0\) 对称端同态 \(A_t\)，使
	\(g_t(v,w)=g_0(A_tv,w)\)，并令 \(B_t=A_t^{-1/2}\)。正定矩阵的正平方根光滑依赖其系数：例如线性化方程 \(C^2=A\) 得 \(CH+HC=K\)，在正定 \(C\) 的特征基中以正数 \(\lambda_i+\lambda_j\) 相除即可解出 \(H\)，有限维隐函数定理遂给出光滑性。因此 \(B_t\) 光滑，且
	\(g_t(B_tv,B_tw)=g_0(v,w)\)。它把 \(g_0\) 的定向正交标架送到 \(g_t\) 的定向正交标架。
	若 \(Q,\eta_0\) 表示原自旋结构，令 \(\eta_t=B_t\circ\eta_0\)，便在同一个主丛 \(Q\) 上得到了运输后的自旋结构。于是所有关联丛可统一写成 \(Q\times_\kappa\Sigma_{2m}\)，手性分次也固定。其 Clifford 作用满足 \(c_t(v)=c_0(B_t^{-1}v)\)，提升联络由式\eqref{g3:eq:spin-connection-local}光滑依赖于 \(t\)。
	用 \(g_0\) 固定 \(H^1(S^+)\) 与 \(L^2(S^-)\) 的范数。紧性保证其他 \(g_t\) 定义的范数与它们一致等价。局部表示中 \(D_{g_t}^+\) 的一阶及零阶系数连续依赖 \(t\)，故
	\(t\mapsto D_{g_t}^+\) 在 \(\mathcal L(H^1(S^+),L^2(S^-))\) 中连续。每个算子都椭圆，\cref{g3:thm:index-local-constancy}遂说明其指标恒定。若某个度量具有正标量曲率，上一推论使它的指标为零，与非零恒定指标矛盾。
\end{proof}

\Cref{g3:thm:spin-ahat-index}将把这个解析指标识别为
\(\langle\widehat A(TM),[M]\rangle\)。在接受那里明确引入的指标定理后，就得到拓扑表述：非零 \(\widehat A\) 数阻碍闭自旋流形上的正标量曲率。奇数维的自旋 Dirac 算子仍满足核消失结论，但本节的手性指标只定义在偶数维中。
% 纲要条目排为编号小节标题，并进入目录与书签；说明段仍保留于纲要注释中。

% !TeX root = ../../Book3.tex
% 纲要标题：### 17.4* Atiyah–Singer 指标定理及其应用
% 纲要位置：几何学计划纲要.md，第 2807 行。
% 正文状态：初稿；尚未正式审核。

% 纲要内容开始
% 
% * 椭圆符号与解析指标
% * 拓扑 K 类和推前
% * Atiyah–Singer 指标定理
% * de Rham、符号差与 spin 指标
% 
% ---
% 
% 纲要内容结束

\OptionalSection{Atiyah–Singer 指标定理及其应用}{g3:sec:1704}

一个椭圆算子的核与余核可以随系数变化，但二者的维数之差保持得更稳定。第十三章已从泛函分析证明这种稳定性；本节把决定这个整数的主符号送入拓扑 \(K\) 理论，再陈述它与拓扑推前相等的指标定理。最后比较三种算子：它们都能写成一个奇的自伴算子的非对角块，所用的分次却不同，因而测量不同的拓扑量。
\ProseParagraphBreak
全节的流形均非空、光滑、紧致且无边界，维数 \(n\ge1\)。向量丛取光滑有限秩复丛，Hermitian 内积第一变量线性。一般指标定理本身采用明确的外部结果；其符号类和推前沿用第二卷第24.4节的构造。本节会证明从本卷 Fredholm、Hodge 与谱理论直接得到的解析识别。

\subsection{椭圆符号与解析指标}
\label{g3:subsec:1704:01}
设 \(P:C^\infty(M;E)\to C^\infty(M;F)\) 是 \(d\ge1\) 阶线性微分算子。在局部标架中写作
\[
	P=\sum_{|\alpha|\le d}a_\alpha(x)\partial^\alpha,
	\qquad
	\sigma_P(x,\xi)=\sum_{|\alpha|=d}a_\alpha(x)(i\xi)^\alpha.
\]
我们继续使用 Fourier 符号约定 \(\partial_j\mapsto i\xi_j\)。最高阶部分的坐标变换规律说明，\(\sigma_P\) 是 \(\pi^*E\to\pi^*F\) 的丛映射，其中 \(\pi:\mathrm T^*M\to M\)；它对纤维变量正齐次 \(d\) 次。椭圆性是 \(\xi\ne0\) 时此映射可逆，特别地，\(E,F\) 在每个分支上的秩相同。

\begin{definition}\label{g3:def:analytic-index}
	闭流形上复丛之间的椭圆算子 \(P\) 的\term[jiexi-zhibiao]{解析指标}[analytic index]定义为
	\begin{equation}\label{g3:eq:analytic-index}
		\operatorname{ind}_{\mathrm a}P
		=\dim_{\mathbb C}\ker P-\dim_{\mathbb C}\ker P^*.
	\end{equation}
	这里 \(P^*\) 是所选密度和丛度量的形式伴随，两个核均由光滑截面组成。
\end{definition}
\symidx{\(\operatorname{ind}_{\mathrm a}P\)：复椭圆算子的解析指标}

\cref{g3:thm:elliptic-fredholm}证明 \(P:H^d(E)\to L^2(F)\) 的像闭且等于 \((\ker P^*)^\perp\)，所以式\eqref{g3:eq:analytic-index}就是其 Fredholm 指标。\cref{g3:cor:elliptic-range-hk}还说明，对每个整数 \(k\ge0\)，改用 \(P:H^{k+d}(E)\to H^k(F)\) 得到同一个整数。更换度量只改变等价范数及余核的正交表示，不改变商空间 \(H^k(F)/P(H^{k+d}(E))\)，因而也不改变指标。

\begin{proposition}\label{g3:prop:index-algebra}
	闭流形上的椭圆算子满足
	\[
		\operatorname{ind}_{\mathrm a}(P\oplus Q)
		=\operatorname{ind}_{\mathrm a}P+\operatorname{ind}_{\mathrm a}Q,
		\qquad
		\operatorname{ind}_{\mathrm a}P^*
		=-\operatorname{ind}_{\mathrm a}P.
	\]
	固定阶数的椭圆系数连续族具有局部不变的指标；给算子加入较低阶项也不改变指标。
\end{proposition}
\begin{proof}
	直和的核与伴随核分别直和，得到第一式。\cref{g3:prop:elliptic-adjoint-domain}把 Hilbert 伴随识别为定义域 \(H^d\) 上的形式伴随；再次取伴随恢复 \(P\)，故第二式成立。固定阶数的系数连续性给出 \(H^d\to L^2\) 的算子范数连续性，应用\cref{g3:thm:index-local-constancy}即可。若 \(B\) 的阶数小于 \(d\)，则 \(P+tB\) 的主符号始终等于 \(\sigma_P\)，沿 \(0\le t\le1\) 使用同一结论。
\end{proof}

形式自伴算子的指标为零，因此几何上有意义的非零指标往往来自它的一个分块。设 \(E=E^+\oplus E^-\) 为正交分次丛，\(\Gamma\) 在两块上分别等于 \(+I,-I\)。一个形式自伴且满足 \(D\Gamma=-\Gamma D\) 的一阶椭圆算子具有形式
\begin{equation}\label{g3:eq:odd-elliptic-block}
	D=\begin{pmatrix}0&D^-\\D^+&0\end{pmatrix},
	\qquad D^-=(D^+)^*,
	\qquad
	D^2=\begin{pmatrix}D^-D^+&0\\0&D^+D^-\end{pmatrix}.
\end{equation}
它的两个非对角块都是椭圆算子。我们要计算的是
\(\operatorname{ind}D^+=\dim\ker D^+-\dim\ker D^-\)，而非总算子 \(D\) 的指标。选定分次也就选定了这个差的顺序。

\subsection{拓扑 K 类和推前}
\label{g3:subsec:1704:02}
先恢复符号类的具体含义。紧 Hausdorff 空间 \(X\) 上的 \(K^0(X)\) 是由复向量丛的形式差生成，并满足 \([E\oplus F]=[E]+[F]\) 的群。对有限 CW 对 \((X,A)\)，相对群可以用三元组 \((E,F,\alpha)\) 描述，其中 \(\alpha:E|_A\to F|_A\) 为同构；允许稳定增加 \((G,G,I)\)、丛同构以及经过可逆映射的同伦。能将 \(\alpha\) 延拓成整个 \(X\) 上同构的三元组代表零。逆元是 \((F,E,\alpha^{-1})\)。第二卷第24.2.1小节已证明，这个模型等于 \(\widetilde K^0(X/A)\)。
% 跨卷依赖：g2:prop:k-triple-model。
\ProseParagraphBreak
对局部紧空间 \(Y\)，定义 \(K_c^0(Y)=\widetilde K^0(Y^+)\)，其中 \(Y^+\) 是单点紧化。这里的“紧支”指两张丛在某个紧集之外由给定映射相识别；它并不要求丛本身在无穷远消失。给 \(\mathrm T^*M\) 选度量，其单位盘丛、球丛记为 \(D(\mathrm T^*M),S(\mathrm T^*M)\)。径向紧化给出
\[
	K_c^0(\mathrm T^*M)
	\cong K^0\bigl(D(\mathrm T^*M),S(\mathrm T^*M)\bigr).
\]

\begin{definition}\label{g3:def:elliptic-symbol-class}
	设 \(M\) 为闭光滑流形，\(E,F\to M\) 为同秩有限秩光滑复向量丛，\(P:C^\infty(E)\to C^\infty(F)\) 为正阶光滑椭圆算子，\(\pi:\mathrm T^*M\to M\) 为投影。选底度量，\(S(\mathrm T^*M)\) 为单位余球丛，\(\sigma_P\) 为主符号。其\term[fuhao-K-lei]{符号 \(K\) 类}[symbol \(K\)-class]是
	\begin{equation}\label{g3:eq:elliptic-symbol-class}
		[\sigma_P]
		=\bigl[\pi^*E,\pi^*F,\sigma_P|_{S(\mathrm T^*M)}\bigr]
		\in K_c^0(\mathrm T^*M).
	\end{equation}
\end{definition}
\symidx{\([\sigma_P]\in K_c^0(\mathrm T^*M)\)：椭圆算子的符号 \(K\) 类}

椭圆性恰好保证第三项在球丛上可逆。改变单位球丛只需沿正径向伸缩，齐次性给出一条可逆同伦，故此类不依赖辅助度量。它也不依赖较低阶项，并随符号的可逆同伦保持不变。直和对应于 \(K\) 群加法。
\ProseParagraphBreak
要从这个类得到整数，还需要推前的定向约定。设 \(V\to X\) 为复秩 \(r\) 丛，\(p:V\to X\)，\(v\) 是 \(p^*V\) 的典范截面。外乘构成复形
\[
	0\longrightarrow p^*\Lambda^0V
	\xrightarrow{\,v\wedge\,}p^*\Lambda^1V
	\xrightarrow{\,v\wedge\,}\cdots
	\xrightarrow{\,v\wedge\,}p^*\Lambda^rV\longrightarrow0.
\]
若 \(v\ne0\)，由度量给出的收缩 \(\iota_v\) 满足
\((v\wedge)\iota_v+\iota_v(v\wedge)=|v|^2I\)，所以这个复形在零截面外正合。把奇偶次数合并，映射
\[
	\varepsilon(v)+\iota_v:
	p^*\Lambda^{\mathrm{ev}}V\longrightarrow p^*\Lambda^{\mathrm{odd}}V
\]
在零截面外可逆，定义 \(K\) 理论的 Thom 类\footnote{René Thom（1923--2002），法国数学家，研究示性类、配边与微分拓扑，后来发展了突变理论。}
\(\lambda_V\)。这里采用的是 \(\Lambda^*V\)，不是 \(\Lambda^*V^*\)。第二卷第24.4.3小节已证明复丛 Thom 同构
\begin{equation}\label{g3:eq:k-thom-index-convention}
	K^0(X)\longrightarrow K_c^0(V),
	\qquad a\longmapsto p^*a\cdot\lambda_V
\end{equation}
以及它的紧支版本；直和满足 \(\lambda_{V\oplus W}=\lambda_V\lambda_W\)。
% 跨卷依赖：g2:prop:complex-k-thom。
\ProseParagraphBreak
这个约定的符号可以在复直线上查清。若 \(H\) 是 \(\mathbb CP^1\) 上典范线丛的对偶，\(\beta=[H]-1\) 是以 \(\int_{\mathbb CP^1}c_1(H)=1\) 归一化的 Bott 元素，\footnote{Raoul Bott（1923--2005），生于匈牙利、长期任教于美国的数学家，研究拓扑与微分几何，建立了拓扑 \(K\) 理论的周期性定理。}则上述外乘复形给出
\begin{equation}\label{g3:eq:index-bott-normalization}
	\lambda_{\mathbb C}=-\beta,
	\qquad
	\lambda_{\mathbb C^N}=(-1)^N\beta^{\boxtimes N}.
\end{equation}
第二卷第24.4.3小节已用圆盘边界的粘接方向证明第一式，第二式由外积得到。定义
\[
	b_N:K_c^0(\mathbb C^N)\xrightarrow{\cong}\mathbb Z,
	\qquad b_N(\lambda_{\mathbb C^N})=1.
\]
以后推前到一点时始终使用这个 \(b_N\)；若把生成元换成 \(\beta^{\boxtimes N}\)，末尾的归一化也必须同时改变。
% 跨卷依赖：g2:eq:k-koszul-bott-sign。

\begin{definition}\label{g3:def:topological-index}
	设 \(M\) 为闭光滑流形。选光滑嵌入 \(i:M\hookrightarrow\mathbb R^N\)，以度量识别 \(\mathrm T^*M\cong\mathrm TM\)。令
	\[
		j=\mathrm Ti:\mathrm TM\hookrightarrow
		\mathrm T\mathbb R^N\cong\mathbb C^N,
		\qquad (x,v)\longmapsto x+iv.
	\]
	用 \(j\) 的复法丛的 Thom 同构及管状邻域延拓定义 \(j_!\)，令 \(b_N:K_c^0(\mathbb C^N)\to\mathbb Z\) 为将标准复 Bott 类送到一的 Bott 同构。\term[tuopu-zhibiao]{拓扑指标}[topological index]为
	\begin{equation}\label{g3:eq:topological-index}
		\operatorname{ind}_{\mathrm t}
		=b_N\circ j_!:
		K_c^0(\mathrm T^*M)\longrightarrow\mathbb Z.
	\end{equation}
\end{definition}
\symidx{\(\operatorname{ind}_{\mathrm t}\)：紧支 \(K\) 理论上的拓扑指标}

这里有两步几何数据需要说明。若 \(\nu\to M\) 是 \(i\) 的实法丛，给切丛选联络后，\(j\) 的法丛识别为
\[
	\pi^*(\nu\oplus\nu)\cong\pi^*(\nu\otimes_{\mathbb R}\mathbb C).
\]
第一份 \(\nu\) 描述底点的法向位移，第二份描述切向量的法向分量；复结构采用 \((a,b)\mapsto a+ib\)。给 \(a\in K_c^0(\mathrm TM)\) 乘这个复法丛的 \(\lambda\) 类，再经管状邻域送入 \(\mathbb C^N\)，最后在邻域外延拓为零，便得到 \(j_!a\)。延拓为零在单点紧化上是压缩邻域补集所诱导的拉回，故不需要把某张非零秩丛直接“补成零丛”。
\ProseParagraphBreak
这里使用的拓扑输入有明确来源：闭光滑流形具有有限 CW 型，见第二卷第17.2.4小节的 Morse 胞腔结论；光滑嵌入及管状邻域的存在见第二卷第15.1、15.3节；相对 \(K\) 群、复丛 Thom 同构及其外积相容性见第二卷第24.2--24.4节。第二卷第24.4.3小节已逐项建立上述推前的选择无关性。简要回顾如下：度量、联络及管状邻域的连续变化给出同伦；把环境增大一个实维数，会在法丛中增加一条复直线，因此 \(j_!a\) 乘上 \(\lambda_{\mathbb C}\)，而 \(b_{N+1}\) 恰好将这个因子消去。任意两个嵌入 \(i_0,i_1\) 在共同环境中由
\[
	x\longmapsto(\cos t\,i_0(x),\sin t\,i_1(x)),
	\qquad 0\le t\le\pi/2,
\]
连接为嵌入族；配合稳定化就得到同一个整数。这也解释了式\eqref{g3:eq:index-bott-normalization}为何关系到推前的良定义。
% 跨卷依赖：g2:cor:morse-cell-type；g2:thm:weak-whitney-compact；g2:thm:tubular-neighborhood；g2:eq:topological-index-pushforward。

\subsection{Atiyah–Singer 指标定理}
\label{g3:subsec:1704:03}
现在两边都有了独立定义：左边是两个有限维核之差，右边由主符号、复法丛和拓扑推前构成。Atiyah--Singer 指标定理把它们识别起来。\footnote{Michael Francis Atiyah（1929--2019），英国数学家，参与建立拓扑 \(K\) 理论及椭圆算子的指标理论。}

\termidx[Atiyah Singer zhibiao dingli]{Atiyah--Singer 指标定理}
\foreignidx{Atiyah--Singer index theorem}
\begin{theorem}[Atiyah--Singer 指标定理]\label{g3:thm:atiyah-singer-index}
	设 \(M\) 为闭光滑流形，\(\dim M\ge1\)，\(E,F\) 为光滑有限秩复向量丛，\(P:C^\infty(M;E)\to C^\infty(M;F)\) 为光滑系数的正阶椭圆微分算子。采用式\eqref{g3:eq:elliptic-symbol-class}的符号类及式\eqref{g3:eq:topological-index}的推前归一化，则
	\begin{equation}\label{g3:eq:atiyah-singer-index}
		\operatorname{ind}_{\mathrm a}P
		=\operatorname{ind}_{\mathrm t}[\sigma_P].
	\end{equation}
	此陈述不要求 \(M\) 定向，也不要求 \(M\) 具有复结构或 spin 结构。
\end{theorem}

这里采用原论文第一篇第6节定理(6.7)在平凡群情形的结论，见\cite[第520页]{AtiyahSinger1968I}。其推前使用第3节的复外代数类，构造和嵌入无关性的讨论见原刊第493--498页，正是上面恢复的约定。完整证明还要建立解析指标对嵌入推前的相容性；本节将这一步作为外部定理使用。
\ProseParagraphBreak
该等式比“加入低阶项不改变指标”更强：它说明所有具有相同紧支 \(K\) 类的符号得到同一指标，即使它们原先作用在不同的丛之间，或需要稳定增加丛之后才相互比较。仅凭第十三章固定 Banach 空间中的局部不变性，不能直接完成这些比较。

\begin{proposition}\label{g3:prop:graded-heat-index}
	设 \((M,g)\) 为非空光滑闭黎曼流形，\(\dim M=n\ge1\)，其上光滑 Hermitian 分次丛 \(E=E^+\oplus E^-\) 带形式自伴奇椭圆算子 \(D\)，令 \(D^+=D|_{E^+}:C^\infty(E^+)\to C^\infty(E^-)\)、\(D^-=D|_{E^-}=(D^+)^*\)，\(\Gamma|_{E^\pm}=\pm I\)；设 \(D\) 为一阶且 \(D^2\) 为 Laplace 型。令 \(K_{D^2}(t,x,y)\) 为其热核，并定义纤维\term[chaoji]{超迹}[supertrace]
	\(\operatorname{str}A=\operatorname{tr}(\Gamma A)\)。则每个 \(t>0\) 都有
	\begin{equation}\label{g3:eq:graded-heat-index}
		\operatorname{ind}D^+
		=\operatorname{Tr}(e^{-tD^-D^+})
		-\operatorname{Tr}(e^{-tD^+D^-})
		=\int_M\operatorname{str}K_{D^2}(t,x,x)\,\mu_g.
	\end{equation}
\end{proposition}
\begin{proof}
	两个对角块非负、形式自伴且椭圆，第十三章的自伴实现和离散谱结论适用。若
	\(D^-D^+u=\lambda u\)，\(\lambda>0\)，则 \(u\) 光滑，且
	\[
		D^+D^-(D^+u)=\lambda D^+u,
		\qquad
		\lVert D^+u\rVert_{L^2}^2=\lambda\lVert u\rVert_{L^2}^2.
	\]
	所以 \(D^+\) 将正特征空间单射到另一块的同一特征空间。反向映射为 \(\lambda^{-1}D^-\)，故二者是同构，正特征值的重数逐一相同。零特征空间分别为 \(\ker D^+\) 与 \(\ker D^-\)。

	\cref{g3:prop:heat-kernel-trace}的丛版本保证两边各自的热迹有限，谱和绝对收敛。因此逐个抵消正谱不会改变和，剩下的恰为两个零特征空间的维数之差。最后对两块分别使用核迹公式，再相减，得到纤维超迹的积分。
\end{proof}

这一步完全来自谱配对，还没有计算任何曲率多项式。若 \(n=2m\)，第十六章的丛热核展开给出
\[
	K_{D^2}(t,x,x)\sim(4\pi t)^{-m}
	\sum_{j\ge0}t^j u_j(x,x).
\]
由于余项在 \(M\) 上统一，式\eqref{g3:eq:graded-heat-index}已经推出
\begin{equation}\label{g3:eq:integrated-supertrace-coefficients}
	\int_M\operatorname{str}u_j(x,x)\,\mu_g=0\quad(j<m),
	\qquad
	\operatorname{ind}D^+
	=(4\pi)^{-m}\int_M\operatorname{str}u_m(x,x)\,\mu_g.
\end{equation}
的确，将统一展开积分，依次乘以最低的 \(t\) 幂并令 \(t\downarrow0\)，即可逐项消去负幂系数，再取常数项。
\ProseParagraphBreak
一般局部指标公式还需要更精细的计算：对相应的几何算子，证明低阶系数的超迹在每一点消失，并将临界阶 \((4\pi)^{-m}\operatorname{str}u_m(x,x)\mu_g\) 识别为指定的示性形式。这涉及曲率在分次模上的作用和临界阶系数，单有 \(u_0\) 及 \(u_1\) 的公式不能处理任意维数。局部指标定理及其热核系数表述见\cite[第4章，4.1节]{BerlineGetzlerVergne2004}。本章前两节以另一条完整的转移形式论证得到 Euler 积分公式；下面的符号差与 spin 示性数公式采用注明的外部结论。

\subsection{de Rham、符号差与 spin 指标}
\label{g3:subsec:1704:04}
先取复外形式丛 \(E=\Lambda^*\mathrm T^*M\otimes\mathbb C\)，按次数奇偶分次，并令 \(D=\mathrm d+\delta\)。由
\[
	\sigma_D(x,\xi)=i\bigl(\varepsilon(\xi)-\iota_{\xi^\sharp}\bigr),
	\qquad \sigma_D(x,\xi)^2=|\xi|^2I,
\]
它是奇的形式自伴椭圆算子，且 \(D^2=\Delta\)。这里不需要定向，因为形式内积与 \(\delta\) 可以直接用黎曼密度定义。

\begin{proposition}\label{g3:prop:de-rham-analytic-index}
	任意闭光滑流形 \(M\) 上，按次数奇偶分次的算子
	\[
		D_{\mathrm{dR}}^+:
		\Omega^{\mathrm{ev}}(M;\mathbb C)
		\longrightarrow\Omega^{\mathrm{odd}}(M;\mathbb C),
		\qquad D_{\mathrm{dR}}^+=(\mathrm d+\delta)|_{\mathrm{ev}},
	\]
	满足 \(\operatorname{ind}D_{\mathrm{dR}}^+=\chi(M)\)。若 \(M\) 定向且维数为 \(2m\)，则
	\[
		\operatorname{ind}D_{\mathrm{dR}}^+
		=\int_M e(\mathrm TM,\nabla^{\mathrm{LC}}).
	\]
\end{proposition}
\begin{proof}
	由\cref{g3:prop:form-hodge-energy}，\(\ker D=\ker\Delta\)。于是
	\[
		\operatorname{ind}D_{\mathrm{dR}}^+
		=\sum_{q=0}^n(-1)^q\dim_{\mathbb C}\mathcal H^q(M)
		=\sum_{q=0}^n(-1)^q\dim_{\mathbb C}H_{\mathrm{dR}}^q(M;\mathbb C),
	\]
	第二个等号来自\cref{g3:thm:hodge-de-rham}。第二卷第7.3.4小节的 de Rham 比较把这些空间识别为奇异上同调；第3.4节的同调 Euler 公式将其交错维数之和识别为 \(\chi(M)\)。最后一式正是本章已证的\cref{g3:thm:gauss-bonnet-chern}。
	% 跨卷依赖：g2:thm:de-rham；g2:thm:euler-homology。
\end{proof}

signature 模型仍使用全部外形式，但换了一个分次。设 \(M\) 闭、定向且 \(\dim M=4k\)，\(k\ge1\)。在 \(q\) 形式上定义
\begin{equation}\label{g3:eq:signature-grading}
	\tau\alpha=(-1)^{q(q-1)/2+k}*\alpha.
\end{equation}
这里 \(*\) 复线性延拓，\(\tau\) 将 \(q\) 次形式送到 \(4k-q\) 次。由
\(*^2=(-1)^{q(4k-q)}\) 计算得 \(\tau^2=I\)；它还是等距映射，故自伴。其 \(\pm1\) 特征丛记为 \(E_\tau^\pm\)。
\ProseParagraphBreak
用\cref{g3:prop:hodge-star-commutes}中的
\(\mathrm d*\alpha=(-1)^q*\delta\alpha\)、
\(\delta*\alpha=(-1)^{q+1}*\mathrm d\alpha\)，并比较相邻次数的因子，就得到
\[
	\tau\mathrm d=-\delta\tau,
	\qquad \tau\delta=-\mathrm d\tau,
	\qquad \tau D=-D\tau.
\]
例如 \(\tau\mathrm d\alpha\) 的因子为
\((-1)^{(q+1)q/2+k}\)，而 \(\delta\tau\alpha\) 的因子为
\((-1)^{q(q-1)/2+k+q+1}\)，二者相差一个负号。因此
\[
	B^+=(\mathrm d+\delta)|_{E_\tau^+}:
	C^\infty(M;E_\tau^+)\longrightarrow C^\infty(M;E_\tau^-)
\]
是一个椭圆算子，称为\term[signature-suanzi]{signature 算子}[signature operator]的正半部分。

\begin{proposition}\label{g3:prop:signature-analytic-index}
	闭定向 \(4k\) 维流形上，中间维实上同调的对称非退化形式
	\[
		Q([\alpha],[\beta])=\int_M\alpha\wedge\beta,
		\qquad [\alpha],[\beta]\in H_{\mathrm{dR}}^{2k}(M;\mathbb R),
	\]
	具有\term[fuhao-cha]{符号差}[signature] \(\operatorname{sign}(M)\)，即其正、负惯性指数之差。任选光滑黎曼度量，以 \(*\) 为复线性 Hodge 星、\(\delta=\mathrm d^*\)，在 \(q\) 次形式上置 \(\tau=(-1)^{q(q-1)/2+k}*\)，并令 \(B^+=(\mathrm d+\delta)|_{\ker(\tau-I)}\) 取值于 \(\ker(\tau+I)\)。则，
	\[
		\operatorname{ind}B^+=\operatorname{sign}(M).
	\]
\end{proposition}
\begin{proof}
	非退化性由\cref{g3:thm:hodge-poincare-duality}得到。由 Hodge 定理，可在实调和形式空间计算 \(Q\)。因为 \(\tau\) 与 \(\Delta=D^2\) 交换，
	\(\operatorname{ind}B^+\) 等于 \(\tau\) 在全部复调和形式上的迹。

	若 \(q<2k\)，则 \(\tau\) 交换 \(\mathcal H^q\) 与 \(\mathcal H^{4k-q}\)。这一对空间上的正负特征空间分别由 \(h+\tau h\) 与 \(h-\tau h\) 参数化，维数相等，故对迹的贡献为零。只剩下中间维。此时
	\[
		(-1)^{(2k)(2k-1)/2+k}=1,
	\]
	所以 \(\tau=*\)。在实中间维调和空间上，\(*^2=I\)，其正负特征空间正交，且
	\[
		Q(\alpha,\beta)=(\alpha,*\beta)_{L^2}.
	\]
	因此 \(Q\) 在正特征空间上正定，在负特征空间上负定，两者互相正交。其符号差就是这两个实维数之差。复化后的复维数等于原实维数，恰好给出 \(\operatorname{ind}B^+\)。
\end{proof}

为写下另外两个示性数，恢复第二卷第21.3.5小节的 Pontryagin 类\footnote{Lev Semenovich Pontryagin（1908--1988），俄罗斯数学家，研究拓扑、拓扑群与控制理论；实向量丛的这组示性类以他命名。}
\(p_j(V)=(-1)^j c_{2j}(V\otimes_{\mathbb R}\mathbb C)\)。
在有理上同调中用形式根记作 \(p(V)=\prod_j(1+x_j^2)\)，其中 \(x_j\) 的形式次数为 \(2\)。定义
\begin{equation}\label{g3:eq:l-ahat-series}
	L(V)=\prod_j\frac{x_j}{\tanh x_j},
	\qquad
	\widehat A(V)=\prod_j\frac{x_j/2}{\sinh(x_j/2)}.
\end{equation}
\symidx{\(L(V),\widehat A(V)\)：实向量丛的 Hirzebruch \(L\) 类与 \(\widehat A\) 类}
每个因子都是 \(x_j\) 的偶幂级数，乘积的齐次部分是 \(x_j^2\) 的对称多项式，因而是 \(p_1,p_2,\ldots\) 的唯一有理多项式。这里无需假设 \(V\) 在原底空间上实际分裂。低次项为
\begin{align}
	L(V)&=1+\frac{p_1}{3}
	+\frac{7p_2-p_1^2}{45}+\cdots,\label{g3:eq:l-low-degrees}\\
	\widehat A(V)&=1-\frac{p_1}{24}
	+\frac{7p_1^2-4p_2}{5760}+\cdots.\label{g3:eq:ahat-low-degrees}
\end{align}
这些类的下标若写成 \(L_{4j},\widehat A_{4j}\)，表示上同调次数 \(4j\)。与基本类配对时只取顶次部分。第二卷第22.3节还给出了相应的曲率代表，因此配对也可以写成微分形式的积分。
% 跨卷依赖：g2:def:pontryagin-classes；g2:eq:pontryagin-curvature-forms。
\ProseParagraphBreak
符号差的拓扑公式是 Hirzebruch 符号差定理。\footnote{Friedrich Hirzebruch（1927--2012），德国数学家，长期任教于波恩大学，研究示性类、符号差与代数几何中的指标公式。}这里将它作为一般指标理论的一个明确的外部几何公式使用。
\termidx[Hirzebruch fuhaocha dingli]{Hirzebruch 符号差定理}
\foreignidx{Hirzebruch signature theorem}
\begin{theorem}\label{g3:thm:signature-l-index}
	设 \(M\) 为闭定向 \(4k\) 维光滑流形，\(k\ge1\)，任选光滑黎曼度量，以 \(*\) 为复线性 Hodge 星、\(\delta=\mathrm d^*\)，在 \(q\) 次形式上置 \(\tau=(-1)^{q(q-1)/2+k}*\)，并令 \(B^+=(\mathrm d+\delta)|_{\ker(\tau-I)}\) 取值于 \(\ker(\tau+I)\)。以 \(L(\mathrm TM)\) 表示由特征级数 \(x/\tanh x\) 定义的 Hirzebruch \(L\) 类，配对取次数 \(4k\) 的部分，则
	\begin{equation}\label{g3:eq:signature-l-index}
		\operatorname{ind}B^+
		=\operatorname{sign}(M)
		=\langle L(\mathrm TM),[M]\rangle.
	\end{equation}
\end{theorem}
该示性数等式及所用 \(L\) 级数可见\cite[第6节，式(6.5)及定理(6.6)，第577页]{AtiyahSinger1968III}。其中解析指标等于符号差的部分，我们已经在上一命题中证明。
\ProseParagraphBreak
最后设 \(M^{2m}\) 带 spin 结构。采用上一节的复自旋子丛 \(S=S^+\oplus S^-\) 及
\[
	\Gamma=i^m c(e_1)\cdots c(e_{2m}),
	\qquad S^\pm=\ker(\Gamma\mp I),
	\qquad D^+:C^\infty(S^+)\to C^\infty(S^-).
\]
\cref{g3:thm:spin-dirac-elliptic}已证明 \(D\) 形式自伴、椭圆并交换自旋子的两个分次；\cref{g3:thm:spin-lichnerowicz}给出其平方。因而 \(\operatorname{ind}D^+\) 是正、负调和自旋子维数之差，并满足式\eqref{g3:eq:graded-heat-index}。

\begin{theorem}\label{g3:thm:spin-ahat-index}
	设 \((M^{2m},g)\) 为闭定向光滑黎曼流形，\(m\ge1\)，并选定一个 spin 结构，\(S^\pm\) 为体积元 \(\Gamma=\mathrm i^m c(e_1)\cdots c(e_{2m})\) 的 \(\pm1\) 特征丛，\(D^+:C^\infty(S^+)\to C^\infty(S^-)\) 为自旋 Dirac 算子；\(\widehat A(\mathrm TM)\) 由特征级数 \((x/2)/\sinh(x/2)\) 定义，配对取次数 \(2m\) 的部分。则
	\begin{equation}\label{g3:eq:spin-ahat-index}
		\operatorname{ind}D^+
		=\langle\widehat A(\mathrm TM),[M]\rangle.
	\end{equation}
	当 \(\dim M\equiv2\pmod4\) 时，右端因没有相应次数的项而为零。
\end{theorem}

该公式在 \(4k\) 维的原始陈述见\cite[第5节，式(5.2)及定理(5.3)，第570--571页]{AtiyahSinger1968III}；同处也指出其余偶数维的无扭曲自旋子指标为零。本节采用的 \(\widehat A\) 归一化明确写在式\eqref{g3:eq:l-ahat-series}中。结合\cref{g3:cor:spin-positive-scalar-vanishing}立即得到：若闭 spin 流形存在处处正标量曲率的度量，则其 \(\widehat A\) 数为零。这个结论中的曲率导致核消失已在上一节证明；将核的指标识别为示性数使用的是刚刚注明的外部定理。

\begin{example}\label{g3:exm:four-dimensional-indices}
	设 \(M\) 为闭定向四维流形。de Rham 指标为 \(\chi(M)\)，signature 指标为
	\(\frac13\langle p_1(\mathrm TM),[M]\rangle\)。若 \(M\) 还带 spin 结构，则
	\[
		\operatorname{ind}D^+
		=-\frac1{24}\langle p_1(\mathrm TM),[M]\rangle
		=-\frac18\operatorname{sign}(M).
	\]
	因此单从这个复指标公式已得到：spin 四维流形的符号差是 \(8\) 的倍数；若还存在正标量曲率度量，则符号差为零。这几个指标并不相同。例如四维球面的调和上同调只在 \(0,4\) 次各有一维，所以 de Rham 指标为 \(2\)，符号差为 \(0\)。
\end{example}

三种模型由同一条解析机制连接：正谱在两个分次之间配对，零谱留下一个整数。次数奇偶分次保留上同调的交错维数；signature 分次在中间维保留交形式的正负差；自旋子分次保留两类调和自旋子的维数差。要把这些整数写成曲率积分，必须进一步识别与各自分次相适应的示性形式。
% 纲要条目排为编号小节标题，并进入目录与书签；说明段仍保留于纲要注释中。


\begin{chaptersummary}
	Euler 形式是曲率矩阵的归一化 Pfaffian。换架公式保证它全局定义，Bianchi 恒等式保证闭性，联络路径给出不同代表之间的转移形式。单位球丛上的典范截面使一个比较联络的 Euler 形式消失，于是拉回的 Euler 形式成为角形式的外微分。

	\ProseParagraphBreak

	角形式在每根球面纤维上的积分为一，这是曲率积分与整数指标连接的归一化。去掉向量场零点后，内边界的负定向和转移公式的负号相消，边界映射度给出各零点指标；Poincaré--Hopf 定理遂把它们的和识别为 Euler 示性数。

	\ProseParagraphBreak

	Dirac 算子的平方把二阶椭圆分析与曲率联系起来。自旋情形的零阶项为 \(\operatorname{Scal}/4\)，从而正标量曲率使调和自旋子消失。一般指标定理进一步把核与余核之差解释为符号的拓扑不变量；de Rham、signature 与 spin 模型分别给出 Euler 示性数、交形式的符号差和 \(\widehat A\) 数，各自的分次、定向与自旋条件决定了公式的适用对象。
\end{chaptersummary}

\BookExercises
\begin{exercise}\label{g3:ex:pfaffian-variation}
	设 \(X(t)\) 为一个开区间上的光滑可逆实 \(2m\) 阶斜对称矩阵族，\(m\ge1\)。证明
	\[
		\frac{\mathrm d}{\mathrm dt}\operatorname{Pf}(X(t))
		=\frac12\operatorname{Pf}(X(t))\operatorname{tr}(X(t)^{-1}\dot X(t)).
	\]
	再用四阶公式说明：Pfaffian 在奇异矩阵处仍有多项式意义下的导数，但题中含逆矩阵的表达式在那里没有定义。
\end{exercise}
\begin{solution}
	先由 \(\det(X+hY)=\det X\det(I+hX^{-1}Y)\) 及行列式的展开得到
	\[
		\left.\frac{\mathrm d}{\mathrm dh}\det(X+hY)\right|_{h=0}
		=\det X\operatorname{tr}(X^{-1}Y).
	\]
	对 \(\operatorname{Pf}(X)^2=\det X\) 求导，得到
	\[
		2\operatorname{Pf}(X)\frac{\mathrm d}{\mathrm dt}\operatorname{Pf}(X)
		=\operatorname{Pf}(X)^2\operatorname{tr}(X^{-1}\dot X).
	\]
	因为 \(X\) 可逆，Pfaffian 非零，可以除去一个 Pfaffian 因子，得到所求公式。这里没有选择平方根的正负支；Pfaffian 本来就是固定的多项式。

	四阶时写上三角变量为 \(a,b,c,d,e,f\)，则
	\[
		\mathrm d\operatorname{Pf}
		=f\,\mathrm da+a\,\mathrm df-e\,\mathrm db-b\,\mathrm de+d\,\mathrm dc+c\,\mathrm dd.
	\]
	例如取 \(X(t)=\left(\begin{smallmatrix}0&t\\-t&0\end{smallmatrix}\right)\oplus\left(\begin{smallmatrix}0&1\\-1&0\end{smallmatrix}\right)\)，则 \(\operatorname{Pf}(X(t))=t\)，在 \(t=0\) 的导数为 \(1\)。此时 \(X(0)^{-1}\) 不存在；多项式导数仍然良好，不能把逆矩阵公式逐项代入这个点。
\end{solution}

\begin{exercise}\label{g3:ex:trivial-plane-nonzero-euler-form}
	令 \(T^2=(\mathbb R/2\pi\mathbb Z)^2\)，取定向 \(\mathrm dx\wedge\mathrm dy\)。在平凡定向 Euclid 平面丛上固定正向正交标架 \((e_1,e_2)\)，令 \(\alpha=\sin x\,\mathrm dy\)，并定义度量联络
	\[
		\nabla e_1=\alpha e_2,\qquad \nabla e_2=-\alpha e_1.
	\]
	计算曲率、Euler 形式及其积分；求从平凡联络到 \(\nabla\) 的转移形式。解释为何这里的 Euler 形式不是恒零的，但它的上同调类为零。
\end{exercise}
\begin{solution}
	\(\mathrm dx,\mathrm dy\) 是从覆盖空间下降的全局一形式，\(\sin x\) 也是周期函数，故联络全局定义。联络矩阵及其曲率为
	\[
		\omega=\begin{pmatrix}0&-\alpha\\\alpha&0\end{pmatrix},\qquad
		\Omega=\begin{pmatrix}0&-\cos x\,\mathrm dx\wedge\mathrm dy\\
		\cos x\,\mathrm dx\wedge\mathrm dy&0\end{pmatrix}.
	\]
	这里 \(\omega\wedge\omega=0\)，因为每项含 \(\alpha\wedge\alpha\)。因此
	\[
		e(\nabla)=-\frac{\cos x}{2\pi}\,\mathrm dx\wedge\mathrm dy,
		\qquad \int_{T^2}e(\nabla)=0.
	\]
	秩二时转移公式没有曲率因子，直接给出
	\[
		\operatorname{TP}(\mathrm d,\nabla)=-\frac{\alpha}{2\pi},\qquad
		\mathrm d\operatorname{TP}(\mathrm d,\nabla)=e(\nabla).
	\]
	该二形式在 \(\cos x\ne0\) 的点非零，却是恰当形式。平凡丛允许平坦联络，也允许题中非平坦联络；联络无关的是 Euler 形式的上同调类，而非每一点的曲率值。
\end{solution}

\begin{exercise}\label{g3:ex:euler-metric-scaling}
	设 \((M^{2m},g)\) 为定向光滑黎曼流形，\(m\ge1\)，\(s>0\) 为常数，\(g_s=s^2g\)。证明两个 Levi-Civita 联络相同，且
	\[
		e(\nabla^{g_s})=e(\nabla^g).
	\]
	若 \(e(\nabla^g)=\mathcal E_g\mathrm{vol}_g\)，求 \(\mathcal E_{g_s}\)。用常曲率公式再次核对这个缩放次数。
\end{exercise}
\begin{solution}
	Koszul 公式的两边同时乘以常数 \(s^2\)，因此确定的联络不变。若 \((e_i)\) 为 \(g\) 的正向正交标架，则 \((s^{-1}e_i)\) 为 \(g_s\) 的正向正交标架，且
	\[
		\nabla(s^{-1}e_j)=\sum_i(s^{-1}e_i)\omega_{ij}.
	\]
	联络矩阵和曲率矩阵在这两组标架中完全相同，代入 Pfaffian 即得 Euler 形式相同。

	体积形式变为 \(\mathrm{vol}_{g_s}=s^{2m}\mathrm{vol}_g\)，所以
	\(\mathcal E_{g_s}=s^{-2m}\mathcal E_g\)。若 \(g\) 的截面曲率为常数 \(\kappa\)，则 \(g_s\) 的截面曲率为 \(s^{-2}\kappa\)。式\eqref{g3:eq:euler-space-form}中的 \(\kappa^m\) 因而乘 \(s^{-2m}\)，与体积因子正好抵消。这项计算同时检查了 Pfaffian 的次数和 \(2m\) 维积分所需的量纲。
\end{solution}

\begin{exercise}\label{g3:ex:euler-product-surfaces}
	设 \((\Sigma_1,g_1)\)、\((\Sigma_2,g_2)\) 为闭连通定向光滑黎曼曲面，其 Gaussian 曲率为 \(K_1,K_2\)。乘积取乘积度量和先第一因子、后第二因子的定向。直接计算
	\[
		e(\nabla^{\Sigma_1\times\Sigma_2})
		=\frac{K_1K_2}{4\pi^2}\,\mathrm{vol}_{g_1}\wedge\mathrm{vol}_{g_2},
	\]
	这里各因子的形式均理解为经投影拉回。再由本章的 Euler 积分公式证明 \(\chi(\Sigma_1\times\Sigma_2)=\chi(\Sigma_1)\chi(\Sigma_2)\)，并比较 \(S^2\times S^2\) 与 \(S^2\times T^2\)。
\end{exercise}
\begin{solution}
	乘积度量的 Koszul 公式表明，联络分别沿两个因子作用，两个因子的切向量之间没有混合曲率。在按因子排列的正向正交标架下，曲率矩阵的两个二阶块分别为
	\[
		\begin{pmatrix}0&K_i\mathrm{vol}_{g_i}\\-K_i\mathrm{vol}_{g_i}&0\end{pmatrix},\qquad i=1,2.
	\]
	因此四阶 Pfaffian 中只有 \(\Omega_{12}\wedge\Omega_{34}\) 保留，给出题中的公式。由 Fubini 定理及\cref{g3:thm:gauss-bonnet-chern}，
	\[
		\begin{aligned}
			\chi(\Sigma_1\times\Sigma_2)
			&=\int_{\Sigma_1\times\Sigma_2}e(\nabla)\\
			&=\left(\frac1{2\pi}\int_{\Sigma_1}K_1\mathrm{vol}_{g_1}\right)
			\left(\frac1{2\pi}\int_{\Sigma_2}K_2\mathrm{vol}_{g_2}\right)
			=\chi(\Sigma_1)\chi(\Sigma_2).
		\end{aligned}
	\]
	两个单位圆球面的乘积有 Euler 示性数 \(4\)，其 Euler 形式为 \(\mathrm{vol}/(4\pi^2)\)。若第二因子取平坦环面，第二块曲率为零，所以整个 Euler 形式为零，\(\chi(S^2\times T^2)=0\)。第一个因子的曲率仍为正；四维 Euler 形式反映的是两个因子曲率的乘积。
\end{solution}

\begin{exercise}\label{g3:ex:rank-two-angular-form}
	设 \(E\to B\) 为定向 Euclid 平面丛，\(\nabla\) 为度量联络。在正向正交标架 \((e_1,e_2)\) 中写
	\[
		\nabla e_1=\alpha e_2,\qquad \nabla e_2=-\alpha e_1,
		\qquad u=\cos\theta\,e_1+\sin\theta\,e_2.
	\]
	由球丛上的转移构造直接证明
	\[
		\Psi_\nabla=\frac{\mathrm d\theta+\alpha}{2\pi}.
	\]
	再把标架旋转角 \(\beta\)，核对这个表达式如何粘合，并验证 \(\mathrm d\Psi_\nabla=-\pi^*e(\nabla)\) 及纤维积分为 \(1\)。这里局部角度只需在适当的小开集上选择。
\end{exercise}
\begin{solution}
	令 \(J\) 为纤维中保持定向的直角旋转，即 \(Je_1=e_2\)、\(Je_2=-e_1\)。在拉回标架中，联络矩阵为 \(\alpha J\)，而
	\[
		Du=(\mathrm d\theta+\alpha)Ju.
	\]
	记 \(q=\mathrm d\theta+\alpha\)。由 \((u,Ju)\) 为正向正交基，修正项
	\(C=u(Du)^{\mathsf T}-(Du)u^{\mathsf T}\)
	满足 \(Cu=-qJu\)、\(CJu=qu\)，所以 \(C=-qJ\)。因此 \(A=D-D^0=qJ\)，其矩阵的 \((1,2)\) 分量为 \(-q\)。秩二时没有曲率因子，式\eqref{g3:eq:sphere-euler-transgression}直接给出 \(\Psi_\nabla=-A_{12}/(2\pi)=q/(2\pi)\)。

	把标架改为
	\[
		e'_1=\cos\beta\,e_1+\sin\beta\,e_2,
		\qquad e'_2=-\sin\beta\,e_1+\cos\beta\,e_2.
	\]
	直接微分得 \(\alpha'=\alpha+\mathrm d\beta\)。同一单位向量的角度变为 \(\theta'=\theta-\beta\)，故
	\(\mathrm d\theta'+\alpha'=\mathrm d\theta+\alpha\)。改变角度分支只加 \(2\pi\) 的整数倍，不改变微分，因此该表达在球丛上确实全局定义。

	曲率矩阵满足 \(\Omega_{12}=-\mathrm d\alpha\)，所以
	\[
		e(\nabla)=-\frac{\mathrm d\alpha}{2\pi},
		\qquad \mathrm d\Psi_\nabla=\frac{\pi^*\mathrm d\alpha}{2\pi}=-\pi^*e(\nabla).
	\]
	式中此前省略的 \(\alpha\) 拉回已显式写出。沿一根正向圆周纤维，\(\alpha\) 的拉回为零，角度 \(\theta\) 增加 \(2\pi\)，因此积分为 \(1\)。
\end{solution}

\begin{exercise}\label{g3:ex:gauss-bonnet-smooth-boundary}
	设 \((M^{2m},g)\) 为紧定向光滑黎曼流形，\(m\ge1\)，允许光滑边界。记 \(\nu\) 为外单位法向，\(\nabla\) 为 Levi-Civita 联络，\(\Psi_\nabla\) 为球丛上的归一化转移形式。证明
	\[
		\int_M e(\nabla)+\int_{\partial M}\nu^*\Psi_\nabla=\chi(M).
	\]
	边界取外法向优先方向。二维时，令 \(J\) 为正向直角旋转，\(T=J\nu\) 为正向单位切向量，并定义有符号测地曲率
	\(k_g=\langle\nabla_TT,JT\rangle\)。证明上式成为
	\[
		\int_M K\,\mathrm dA+\int_{\partial M}k_g\,\mathrm ds=2\pi\chi(M),
	\]
	并用 Euclid 圆盘检查边界项的符号。
\end{exercise}
\begin{solution}
	球丛构造及 \(\mathrm d\Psi_\nabla=-\pi^*e(\nabla)\) 都是光滑系数的局部计算，因此在边界附近把坐标系数光滑延拓后，同样成立到边界。正定性在足够小的延拓邻域中保持；最终公式与延拓选择无关。

	第二卷第17.1.4小节的带边界 Morse 函数构造，给出边界为最高水平、全部临界点在内部且非退化的函数。它对给定度量的梯度 \(X\) 也严格外向：在边界上微分沿切方向为零、沿内向方向为负，故 \(\langle X,\nu\rangle=\mathrm df(\nu)>0\)。这样选取的 \(X\) 只有有限个内部零点。第二卷第17.4.4小节的带边界 Poincaré--Hopf 定理给出
	\[
		\sum_p\operatorname{ind}_p(X)=\chi(M).
	\]
	若有无边界连通分支，在这些分支另选 Morse 梯度场即可。

	挖去各零点周围的不交小球 \(B_i\)，令 \(M_0=M\setminus\bigcup_i\operatorname{Int}B_i\)。它的边界是原有外边界和负向的小球面之并，故
	\[
		-\int_{M_0}e(\nabla)
		=\int_{\partial M}\widehat X^*\Psi_\nabla
		-\sum_i\int_{\partial B_i}\widehat X^*\Psi_\nabla.
	\]
	代入式\eqref{g3:eq:angular-local-index-integral}，将球内 Euler 积分移回左侧，得到
	\[
		\int_Me(\nabla)+\int_{\partial M}\widehat X^*\Psi_\nabla
		=\sum_i\operatorname{ind}_{p_i}(X)=\chi(M).
	\]
	剩下要把边界上的 \(\widehat X\) 换成 \(\nu\)。向量 \((1-t)\widehat X+t\nu\) 的外法向分量始终为正，故归一化给出它们之间的光滑单位截面同伦 \(H:[0,1]\times\partial M\to S(\mathrm TM)\)。其底映射与 \(t\) 无关，所以顶次形式 \((\pi\circ H)^*e(\nabla)\) 为零。再次用 Stokes 定理，
	\[
		\int_{\partial M}\nu^*\Psi_\nabla
		-\int_{\partial M}\widehat X^*\Psi_\nabla
		=\int_{[0,1]\times\partial M}H^*\mathrm d\Psi_\nabla=0.
	\]
	由此得到第一式。

	二维时，\cref{g3:ex:rank-two-angular-form}说明对任意单位截面 \(u\)，其转移形式拉回在向量 \(Y\) 上的值为
	\(\langle\nabla_Yu,Ju\rangle/(2\pi)\)。取 \(u=\nu\)、\(Y=T=J\nu\)，有
	\[
		2\pi(\nu^*\Psi_\nabla)(T)
		=\langle\nabla_T\nu,T\rangle
		=-\langle\nu,\nabla_TT\rangle
		=\langle\nabla_TT,JT\rangle=k_g,
	\]
	因为 \(JT=-\nu\)。再用 \(e(\nabla)=K\,\mathrm dA/(2\pi)\)，即得曲面公式。半径 \(R\) 的 Euclid 圆盘取逆时针边界，\(K=0\)、\(k_g=1/R\)、边界长度为 \(2\pi R\)，故左边为 \(2\pi=2\pi\chi(D^2)\)。
\end{solution}

\begin{exercise}\label{g3:ex:unoriented-euler-density}
	设 \((M^{2m},g)\) 为闭光滑黎曼流形，\(m\ge1\)，不要求可定向。在每个定向坐标域内，令 \(e_{\mathrm{loc}}(\nabla)\) 为按照该方向计算的切丛 Euler 形式，并写
	\[
		e_{\mathrm{loc}}(\nabla)=F_g\,\mathrm{vol}_{g,\mathrm{loc}}.
	\]
	证明 \(F_g\) 拼成全局光滑函数，密度 \(F_g\,\mu_g\) 满足
	\[
		\int_MF_g\,\mu_g=\chi(M).
	\]
	使用定向双覆盖完成证明，并计算单位圆球面商度量下 \(\mathbb{RP}^{2m}\) 的该积分。证明中说明为何双覆盖的 Euler 示性数为原来的两倍。
\end{exercise}
\begin{solution}
	反转局部方向时，切丛 Pfaffian 与局部体积形式都反号，故比值 \(F_g\) 不变。在方向一致的交集上两者本来就分别相同，因此 \(F_g\) 是全局光滑函数。它可以有正负号，\(F_g\,\mu_g\) 是有符号密度，而非必为正的测度。

	令 \(p:\widetilde M\to M\) 为定向双覆盖：其点是在底点处选择的一种切空间方向。它具有典范定向，取拉回度量后 \(p\) 局部等距。联络与曲率在局部等距下自然，故
	\[
		e(\nabla^{\widetilde M})=(F_g\circ p)\,\mathrm{vol}_{\widetilde g}.
	\]
	在均匀覆盖坐标域上每个积分出现两份，用有限单位分解相加，得到
	\[
		\int_{\widetilde M}e(\nabla^{\widetilde M})
		=2\int_MF_g\,\mu_g.
	\]
	这也适用于可定向分支，其定向双覆盖由两个分支组成。

	可直接用已证的向量场指标公式核对 Euler 示性数的倍数。选取 \(M\) 上一个 Morse 函数的梯度场 \(X\)，由 \(\mathrm dp\) 定义提升场 \(\widetilde X\)。每个零点恰有两个提升；在由 \(p\) 对应的局部坐标中，两个场的分量函数相同，因此两个提升零点的局部指标都等于原指标。第二卷第17.4.3小节的 Poincaré--Hopf 定理不要求流形可定向，故
	\[
		\chi(\widetilde M)=\sum_{\widetilde X(q)=0}\operatorname{ind}_q(\widetilde X)
		=2\sum_{X(x)=0}\operatorname{ind}_x(X)=2\chi(M).
	\]
	在 \(\widetilde M\) 上应用\cref{g3:thm:gauss-bonnet-chern}，将前两式相除以 \(2\)，即得所求公式。

	对径映射在 \(S^{2m}\) 上反向，商映射 \(S^{2m}\to\mathbb{RP}^{2m}\) 正是定向双覆盖，并且对题中度量局部等距。球面高度函数只有两极两个临界点，其梯度的指标分别为 \((-1)^{2m}=1\) 与 \(1\)，因此 \(\chi(S^{2m})=2\)。由上面的二倍关系，\(\chi(\mathbb{RP}^{2m})=1\)，Euler 密度积分也为 \(1\)。这项计算不需要给射影空间选一个不存在的整体方向。
\end{solution}

\begin{exercise}\label{g3:ex:angular-integral-arbitrary-index}
	在 \(\mathbb R^{2m}=\mathbb C^m\) 上采用复坐标给出的实方向，\(m\ge1\)。对整数 \(k\ge1\)，考虑平凡定向实秩 \(2m\) 丛的截面
	\[
		F_k(z)=(z_1^k,z_2,\ldots,z_m),\qquad
		F_{-k}(z)=(\overline z_1^{\,k},z_2,\ldots,z_m),
		\qquad F_0(z)=(|z|^2,0,\ldots,0).
	\]
	证明它们只有原点一个零点，局部指标分别为 \(k,-k,0\)。若 \(\nabla\) 为该丛上任意光滑度量联络，求 \(\widehat F_j=F_j/|F_j|\) 在半径 \(\varepsilon\) 球面上的转移积分；再说明平坦联络与一般联络所得积分的区别。归一化使用平凡丛的标准内积。
\end{exercise}
\begin{solution}
	三个映射的零点都只有原点。对 \(F_k\)，在固定小球的边界上有正下界 \(c=\min|F_k|>0\)。选小的正实数 \(a\)，使 \(a<c\)、\(a^{1/k}<\varepsilon\)。同伦
	\[
		F_k(z)-t(a,0,\ldots,0),\qquad 0\le t\le1,
	\]
	始终在边界非零。在 \(t=1\) 时，零点正是 \(z_1^k=a\)、\(z_2=\cdots=z_m=0\) 的 \(k\) 个解。每个解处，第一复坐标的导数是乘以非零复数 \(kz_1^{k-1}\)，其二阶实行列式为 \(|kz_1^{k-1}|^2>0\)；其他坐标的导数为恒等。因此全部零点非退化且各指标为 \(+1\)。第二卷第16.3.3小节的边界同伦不变性与局部指标可加性给出 \(\operatorname{ind}_0(F_k)=k\)。

	对 \(F_{-k}\) 作同样的平移。每个零点处，第一坐标的实微分是一个非零复数乘法与复共轭的复合。复共轭的实行列式为 \(-1\)，所以每个零点的指标为 \(-1\)，总和为 \(-k\)。而 \(F_0/|F_0|\) 在每个小球面上都是常值 \((1,0,\ldots,0)\)，其映射度为零。因此任何整数都可以由这种光滑孤立零点实现。

	将式\eqref{g3:eq:angular-local-index-integral}用于平凡丛和这些截面，得到对每个整数 \(j\)，
	\[
		\int_{\partial B_\varepsilon}\widehat F_j^*\Psi_\nabla
		=j-\int_{B_\varepsilon}e(\nabla).
	\]
	这里球面取标准实方向诱导的边界方向。若联络平坦，Euler 形式恒为零，转移积分对每个半径都恰等于整数 \(j\)。对一般度量联络，球内 Euler 积分可能非零，因此固定半径上的转移积分未必是整数；其与 \(j\) 的差为 \(O(\varepsilon^{2m})\)，在半径趋零时仍趋于 \(j\)。局部指标来自边界映射度，曲率修正则记录所用联络在球内的信息。
\end{solution}

\begin{optionalexercise}\label{g3:ex:spin-two-half-angle}
	在 \(\Sigma_2=\Lambda^*\mathbb C\) 的基 \((1,e)\) 下，写出 \(c(e_1),c(e_2),\Gamma\) 的矩阵，计算 \(s(t)=\cos t+\sin t\,e_1e_2\) 的自旋表示。证明该表示不能下降为 \(\operatorname{SO}(2)\) 的表示，并求平坦二维空间中 \(D\) 的矩阵微分表达式。
\end{optionalexercise}
\begin{solution}
	外乘和收缩的矩阵分别为
	\[
		\varepsilon=\begin{pmatrix}0&0\\1&0\end{pmatrix},\qquad
		\iota=\begin{pmatrix}0&1\\0&0\end{pmatrix}.
	\]
	所以
	\[
		c(e_1)=\begin{pmatrix}0&-1\\1&0\end{pmatrix},\qquad
		c(e_2)=\begin{pmatrix}0&\mathrm i\\\mathrm i&0\end{pmatrix},\qquad
		\Gamma=\mathrm i c(e_1)c(e_2)=\begin{pmatrix}1&0\\0&-1\end{pmatrix}.
	\]
	乘积 \(c(e_1)c(e_2)=\operatorname{diag}(-\mathrm i,\mathrm i)\)，从而
	\[
		\kappa(s(t))=\begin{pmatrix}e^{-\mathrm it}&0\\0&e^{\mathrm it}\end{pmatrix}.
	\]
	\(\rho(s(t))\) 是旋转角 \(2t\) 的正交映射，因此物理空间中的一次 \(2\pi\) 旋转由 \(s(\pi)=-1\) 表示，它在自旋子上作用为 \(-\operatorname{Id}\)。如果存在 \(\widetilde\kappa:\operatorname{SO}(2)\to\operatorname{GL}(\Sigma_2)\) 使 \(\kappa=\widetilde\kappa\circ\rho\)，则 \(\kappa(-1)=\widetilde\kappa(1)=\operatorname{Id}\)，矛盾。
	在标准平坦标架中，联络项为零，故
	\[
		D=\begin{pmatrix}
			0&-\partial_{x_1}+\mathrm i\partial_{x_2}\\
			\partial_{x_1}+\mathrm i\partial_{x_2}&0
		\end{pmatrix}.
	\]
	两个非对角微分算子的乘积均为 \(-\partial_{x_1}^2-\partial_{x_2}^2\)，直接检验了平坦情形的 \(D^2=\Delta\operatorname{Id}\)。
\end{solution}

\begin{optionalexercise}\label{g3:ex:flat-torus-spin-spectrum}
	设 \(T^n=\mathbb R^n/(L_1\mathbb Z\oplus\cdots\oplus L_n\mathbb Z)\)，\(n\ge2\)，各 \(L_j>0\)，取标准平坦度量和定向。对每个 \(\epsilon=(\epsilon_1,\ldots,\epsilon_n)\in\{0,1/2\}^n\)，构造其平移提升的符号为 \(e^{2\pi\mathrm i\epsilon_j}\) 的自旋结构。证明这给出全部 \(2^n\) 个等价类，并计算各结构的 Dirac 谱及零特征空间。对非零频率，应给出正、负特征值各自的重数；不同频率产生同一数值时再累加重数。
\end{optionalexercise}
\begin{solution}
	以 \(\ell\in\mathbb Z^n\) 表示平移 \(x\mapsto x+(L_1\ell_1,\ldots,L_n\ell_n)\)，定义符号特征
	\[
		\chi_\epsilon(\ell)=\exp\left(2\pi\mathrm i\sum_j\epsilon_j\ell_j\right)\in\{1,-1\}.
	\]
	让 \(\mathbb Z^n\) 作用于 \(\mathbb R^n\times\operatorname{Spin}(n)\)：
	\[
		\ell\cdot(x,s)=(x+L\ell,\chi_\epsilon(\ell)s).
	\]
	这个作用自由且适当不连续，并与右 \(\operatorname{Spin}(n)\) 作用交换。商主丛 \(Q_\epsilon\) 通过 \([x,s]\mapsto([x],\rho(s))\) 映到平凡的定向正交标架丛，给出自旋结构。两种不同 \(\epsilon\) 在某个坐标圆周上的提升具有不同的端点符号，因而其差的一上闭链代表不同的 \(H^1(T^n;\mathbb F_2)\) 类。由于坐标圆周给出 \(H_1(T^n;\mathbb Z)=\mathbb Z^n\)，有 \(H^1(T^n;\mathbb F_2)=\mathbb F_2^n\)；第二卷的分类结论说明这些 \(2^n\) 个结构恰为全部。
	对应自旋子可看成满足
	\[
		\psi(x+L_j e_j)=e^{2\pi\mathrm i\epsilon_j}\psi(x)
	\]
	的 \(\Sigma_n\) 值函数。乘以标量因子 \(\exp(-2\pi\mathrm i\sum_j\epsilon_j x_j/L_j)\) 后成为周期函数，所以普通 Fourier 展开给出完整正交分解
	\[
		\psi(x)=\sum_{k\in\mathbb Z^n}e^{2\pi\mathrm i\langle\xi_k,x\rangle}v_k,
		\qquad (\xi_k)_j=\frac{k_j+\epsilon_j}{L_j},\quad v_k\in\Sigma_n.
	\]
	平坦自旋联络在覆盖空间上是普通微分，因此第 \(k\) 个频率块上的算子为
	\[
		D_k=2\pi\mathrm i c(\xi_k),\qquad D_k^2=4\pi^2|\xi_k|^2\operatorname{Id}.
	\]
	若 \(\xi_k\ne0\)，取单位向量 \(a\perp\xi_k\)，这是 \(n\ge2\) 保证的。可逆映射 \(c(a)\) 与 \(\mathrm i c(\xi_k)\) 反交换，故将正、负特征空间互相同构。于是两个特征值是
	\[
		\lambda_{k,\pm}=\pm2\pi\left(\sum_{j=1}^n\frac{(k_j+\epsilon_j)^2}{L_j^2}\right)^{1/2},
	\]
	每个频率块中各有重数 \(2^{\lfloor n/2\rfloor-1}\)。Fourier 基的完备性保证没有遗漏的谱；更具体地，块上单位特征基的并给出 \(L^2\) 的正交特征基，\(H^1\) 定义域恰对应带权平方和有限的系数。
	零频率出现当且仅当所有 \(\epsilon_j=0\)，此时只能是 \(k=0\)，核为所有常值自旋子，维数 \(2^{\lfloor n/2\rfloor}\)。其余自旋结构的核为零。因此同一平坦度量上的自旋结构选择会改变零模数。
\end{solution}

\begin{optionalexercise}\label{g3:ex:spin-sharper-eigenvalue-bound}
	设 \((M^n,g)\) 为非空闭连通定向自旋流形，\(n\ge2\)，且 \(\operatorname{Scal}_g\ge s_0>0\)。对光滑自旋子定义
	\[
		P_X\psi=\nabla^S_X\psi+\frac1n c(X)D\psi.
	\]
	证明任意非零 Dirac 特征自旋子 \(D\psi=\lambda\psi\) 满足
	\[
		\lambda^2\ge\frac{n}{4(n-1)}s_0.
	\]
	证明等号成立当且仅当该特征自旋子满足 \(\nabla^S_X\psi=-(\lambda/n)c(X)\psi\) 且 \(\operatorname{Scal}_g\equiv s_0\)。
\end{optionalexercise}
\begin{solution}
	用任意正交归一标架求和，并利用 \(c(e_j)^*=-c(e_j)\)，得到
	\begin{align*}
		\sum_j|P_{e_j}\psi|^2
		&=|\nabla^S\psi|^2+\frac1n|D\psi|^2
		+\frac2n\operatorname{Re}\sum_j\langle\nabla^S_{e_j}\psi,c(e_j)D\psi\rangle\\
		&=|\nabla^S\psi|^2-\frac1n|D\psi|^2.
	\end{align*}
	将其积分，再代入式\eqref{g3:eq:spin-integrated-lichnerowicz}，便得
	\[
		\frac{n-1}{n}\|D\psi\|_2^2
		=\|P\psi\|_2^2+\frac14\int_M\operatorname{Scal}_g|\psi|^2\,\mu_g
		\ge\frac{s_0}{4}\|\psi\|_2^2.
	\]
	特征值是实数，代入 \(D\psi=\lambda\psi\) 并除以非零范数平方，得到下界。
	若等号成立，则 \(P\psi=0\)，且 \((\operatorname{Scal}_g-s_0)|\psi|^2=0\) 处处成立。前一式给出题述协变导数关系。由于 \(\lambda\) 实且 \(c(X)\) 斜 Hermitian，
	\[
		 X(|\psi|^2)=2\operatorname{Re}\langle\nabla^S_X\psi,\psi\rangle
		=-\frac{2\lambda}{n}\operatorname{Re}\langle c(X)\psi,\psi\rangle=0.
	\]
	连通性及 \(\psi\not\equiv0\) 说明其长度是正的常数，故 \(\operatorname{Scal}_g\equiv s_0\)。反过来，若两条件成立，则上面的积分恒等式右边恰为 \((s_0/4)\|\psi\|_2^2\)，所以得到等号。
\end{solution}

\begin{optionalexercise}\label{g3:ex:forms-dirac-not-spin-dirac}
	在光滑无边界黎曼流形上，给 \(E=\Lambda^*T^*M\otimes\mathbb C\) 配备诱导联络以及 \(c(v)=\varepsilon(v^\flat)-\iota_v\)。证明其 Dirac 算子是 \(\mathrm d+\delta\)，平方为各次数 Hodge Laplace 算子的直和。用单位圆球面 \(S^2\) 上的一个具体截面说明：一般 Clifford 模丛的曲率项不能直接替换成 \(\operatorname{Scal}/4\)。
\end{optionalexercise}
\begin{solution}
	楔乘与收缩的反交换关系给出 Clifford 关系；二者互为伴随，故 \(c(v)\) 斜 Hermitian。诱导联络保持度量、楔乘和收缩，因而满足 Clifford 相容性。取一点的正规正交标架，第十五章的局部公式给出
	\[
		\mathrm d=\sum_j\varepsilon(e^j)\nabla_{e_j},\qquad
		\delta=-\sum_j\iota_{e_j}\nabla_{e_j}.
	\]
	所以 \(D_E=\sum_jc(e_j)\nabla_{e_j}=\mathrm d+\delta\)。由 \(\mathrm d^2=\delta^2=0\)，有 \(D_E^2=\mathrm d\delta+\delta\mathrm d\)，且此算子保留形式次数。
	单位 \(S^2\) 的标量曲率恒为 \(2\)。常数 \(0\)-形式 \(1\) 是一个非零平行截面，并满足 \(D_E1=0\)。若对这个模丛也有 \(D_E^2=\nabla^*\nabla+\operatorname{Scal}/4\)，把两边作用于 \(1\) 就会得到 \(0=1/2\)，矛盾。正确的式\eqref{g3:eq:dirac-square-general-curvature}使用外形式丛的曲率；该曲率在 \(0\)-形式上作用为零，因而与上述计算相符。自旋子丛的曲率公式是 Lichnerowicz 标量项成立的关键条件。
\end{solution}

\begin{optionalexercise}\label{g3:ex:index-bott-stabilization}
	设 \(M\) 为闭流形，\(a\in K_c^0(\mathrm T^*M)\)。对嵌入
	\(i:M\hookrightarrow\mathbb R^N\)，设切映射推前满足
	\[
		(\mathrm Ti)_!a=q_N\beta^{\boxtimes N},
		\qquad q_N\in\mathbb Z,
	\]
	其中 \(\beta=[H]-1\) 是式\eqref{g3:eq:index-bott-normalization}的正生成元。
	证明增加一个零的环境坐标后 \(q_{N+1}=-q_N\)，而
	\((-1)^Nq_N\) 不变。将此应用于 \(S^2\) 上的 de Rham 符号，分别求标准嵌入 \(S^2\hookrightarrow\mathbb R^3\) 及其 \(\mathbb R^4\) 中稳定化所给的 \(q_N\)。说明为什么不能一面沿用 \(\Lambda^*V\) 的 Thom 类，一面把 \(q_N\) 直接叫作拓扑指标。
\end{optionalexercise}
\begin{solution}
	环境从 \(\mathbb R^N\) 增到 \(\mathbb R^{N+1}\) 时，切映射法丛增加一条复直线。Thom 类的乘法公式及
	\(\lambda_{\mathbb C}=-\beta\) 给出
	\[
		(\mathrm Ti_{\mathrm{new}})_!a
		=(\mathrm Ti)_!a\boxtimes\lambda_{\mathbb C}
		=-q_N\beta^{\boxtimes(N+1)}.
	\]
	因此 \(q_{N+1}=-q_N\)。另一方面，
	\[
		b_N\bigl(q_N\beta^{\boxtimes N}\bigr)
		=(-1)^Nq_N,
	\]
	因为 \(b_N(\lambda_{\mathbb C^N})=1\)。增加一个维数恰好同时改变两个负号，故此整数稳定。

	球面 \(S^2\) 的实上同调在 \(0,2\) 次各一维，其余为零。由\cref{g3:prop:de-rham-analytic-index,g3:thm:atiyah-singer-index}，该符号的拓扑指标为 \(2\)。于是
	\[
		(-1)^3q_3=2,\qquad (-1)^4q_4=2,
		\qquad q_3=-2,\quad q_4=2.
	\]
	若直接取 \(q_N\)，同一个符号仅因增加一个无用的环境坐标就改变所谓指标。这不是解析指标的性质，而是把两种生成元归一化混用了。采用 \(\beta^{\boxtimes N}\) 当然也能描述 \(K\) 群，但此时推前到整数的映射必须带上 \((-1)^N\)。
\end{solution}

\begin{optionalexercise}\label{g3:ex:cp4-index-comparison}
	给 \(M=\mathbb CP^4\) 赋复定向。利用第二卷第7.4节与第21.4.1小节已经得到的
	\[
		H^*(M;\mathbb R)=\mathbb R[h]/(h^5),
		\qquad \langle h^4,[M]\rangle=1,
		\qquad p(\mathrm TM)=(1+h^2)^5,
	\]
	其中 \(h\) 是超平面类的实系数像，计算 \(\chi(M)\)、\(\operatorname{sign}(M)\)、\(\langle L_8(\mathrm TM),[M]\rangle\) 及
	\(\langle\widehat A_8(\mathrm TM),[M]\rangle\)。用指标的整性说明 \(M\) 不存在 spin 结构，并区分本题中哪几个数已经是实际定义的算子的指标。
	% 跨卷依赖：g2:eq:projective-de-rham-ring；g2:eq:complex-projective-pontryagin。
\end{optionalexercise}
\begin{solution}
	非零上同调空间在 \(0,2,4,6,8\) 次各为一维，因此 \(\chi(M)=5\)。中间维空间 \(H^4(M;\mathbb R)\) 由 \(h^2\) 生成，且
	\[
		Q(h^2,h^2)=\langle h^4,[M]\rangle=1.
	\]
	所以交形式为一维正定，\(\operatorname{sign}(M)=1\)。

	展开切丛公式，得到
	\[
		p_1=5h^2,\qquad p_2=10h^4,\qquad p_1^2=25h^4.
	\]
	由式\eqref{g3:eq:l-low-degrees}，
	\[
		\langle L_8,[M]\rangle
		=\frac{7\cdot10-25}{45}=1,
	\]
	与直接算出的符号差一致。而式\eqref{g3:eq:ahat-low-degrees}给出
	\[
		\langle\widehat A_8,[M]\rangle
		=\frac{7\cdot25-4\cdot10}{5760}
		=\frac{135}{5760}=\frac3{128}.
	\]
	若 \(M\) 有 spin 结构，\cref{g3:thm:spin-ahat-index}将使这个数等于两个有限维复向量空间维数之差，因而是整数，与 \(3/128\) 矛盾。

	de Rham 算子与signature 算子已经在 \(M\) 上定义，指标分别为 \(5\) 与 \(1\)。\(\widehat A\) 类作为 Pontryagin 类的有理多项式也能在 \(M\) 上定义，但这里不存在无扭曲 spin 自旋子丛，故 \(3/128\) 不能被解释为本节 spin Dirac 算子的指标。
\end{solution}

\begin{optionalexercise}\label{g3:ex:surface-product-signature}
	令 \(\Sigma_g,\Sigma_h\) 为亏格分别为 \(g,h\ge0\) 的闭连通定向曲面，给 \(M=\Sigma_g\times\Sigma_h\) 赋乘积定向。使用域系数的乘积上同调分解，直接求 \(H^2(M;\mathbb R)\) 上交形式的正、负惯性指数，并比较 de Rham 指标与 signature 指标。再利用切丛直和说明 \(p_1(\mathrm TM)=0\)，核对符号差公式。
\end{optionalexercise}
\begin{solution}
	选归一化的顶次类 \(a\in H^2(\Sigma_g;\mathbb R)\)、
	\(b\in H^2(\Sigma_h;\mathbb R)\)，其积分均为 \(1\)。以后将它们与向 \(M\) 的拉回用同一字母表示。选第一张曲面的辛基
	\(u_i,v_i\)（\(1\le i\le g\)）及第二张曲面的辛基
	\(s_j,t_j\)（\(1\le j\le h\)），使
	\[
		\int_{\Sigma_g}u_i\wedge v_{i'}=\delta_{ii'},
		\qquad
		\int_{\Sigma_h}s_j\wedge t_{j'}=\delta_{jj'},
	\]
	同类基向量的配对为零。辛基的存在可以逐次选一对配对非零的向量、归一化，再取其配对正交补，有限归纳得到。

	乘积分解为
	\[
		H^2(M;\mathbb R)
		=\operatorname{span}\{a,b\}
		\oplus
		\bigl(H^1(\Sigma_g;\mathbb R)\otimes H^1(\Sigma_h;\mathbb R)\bigr).
	\]
	第一块上 \(a^2=b^2=0\)、\(\int_Ma\wedge b=1\)，所以矩阵为
	\(\left(\begin{smallmatrix}0&1\\1&0\end{smallmatrix}\right)\)，贡献一正一负。该块与第二块正交，因为相乘后至少有一张曲面上的形式次数超过 \(2\)。

	第二块分解为 \(gh\) 个互相正交的四维空间。固定 \(i,j\)，令
	\[
		A=u_i\wedge s_j,\quad B=u_i\wedge t_j,\quad
		C=v_i\wedge s_j,\quad D=v_i\wedge t_j.
	\]
	将第二张曲面的一形式移过第一张曲面的一形式会产生一个负号。因此该块的非零配对恰为
	\[
		Q(A,D)=-1,\qquad Q(B,C)=1,
	\]
	以及对称位置的相同数值。它是两个双曲平面的直和，贡献两正两负。综上，
	\[
		b_2^+(M)=1+2gh,\qquad b_2^-(M)=1+2gh,
		\qquad\operatorname{sign}(M)=0.
	\]
	若 \(g=0\) 或 \(h=0\)，相应的四维块为空，上述公式仍成立。

	乘积上同调还给出
	\[
		b_0=b_4=1,\qquad b_1=b_3=2g+2h,\qquad b_2=2+4gh.
	\]
	于是 de Rham 指标为
	\(\chi(M)=4(1-g)(1-h)\)，signature 指标为 \(0\)。最后
	\[
		\mathrm TM\cong
		\pi_1^*\mathrm T\Sigma_g\oplus\pi_2^*\mathrm T\Sigma_h.
	\]
	第一 Pontryagin 类对直和相加，而两张曲面上四次上同调均为零，故 \(p_1(\mathrm TM)=0\)。符号差公式
	\(\operatorname{sign}(M)=\frac13\langle p_1,[M]\rangle\)
	再次给出零。这里符号差的消失已经由交形式直接证明，示性类公式提供了独立的比较。
\end{solution}

\begin{optionalexercise}\label{g3:ex:spin-eight-dimensional-numbers}
	设 \(M\) 为闭定向 spin 八维流形，记
	\[
		a=\langle p_1(\mathrm TM)^2,[M]\rangle,\qquad
		b=\langle p_2(\mathrm TM),[M]\rangle.
	\]
	用 \(a,b\) 写出符号差 \(s\) 与 spin Dirac 指标 \(I\)，再反解 \(a,b\)。若 \(M\) 存在正标量曲率度量，\(a,b,s\) 必须满足什么关系？特别地，说明任何满足 \(a=896,b=128\) 的此类流形均不可能存在正标量曲率度量。本题不要求证明具有这组示性数的流形存在。
\end{optionalexercise}
\begin{solution}
	八次齐次项给出
	\[
		s=\frac{7b-a}{45},\qquad
		I=\frac{7a-4b}{5760}.
	\]
	消去 \(a\)，有
	\(45b=5760I+315s\)，从而
	\[
		b=7s+128I,\qquad a=4s+896I.
	\]
	因为 \(s,I\) 均为整数，示性数必须落在由
	\((4,7)\)、\((896,128)\) 生成的整数格中。等价地，
	\(7b-a\) 必须被 \(45\) 整除，\(7a-4b\) 必须被 \(5760\) 整除；这些是符号差公式和 spin 指标整性共同施加的必要条件。

	若存在正标量曲率度量，Lichnerowicz 公式使所有调和自旋子消失，所以 \(I=0\)。于是
	\[
		a=4s,\qquad b=7s,\qquad 7a=4b.
	\]
	这些是必要条件，并未从中推出正标量曲率度量的存在。对题设数值，
	\[
		s=\frac{7\cdot128-896}{45}=0,\qquad
		I=\frac{7\cdot896-4\cdot128}{5760}=1.
	\]
	非零 spin Dirac 指标给出所求障碍。这个例子也说明符号差为零与自旋子指标为零是不同的条件，八维时一个并不由另一个决定。
\end{solution}
